\documentclass[aps,preprint]{revtex4}%
\usepackage{amssymb}
\usepackage{amsfonts}
\usepackage{amsmath}
\usepackage{graphicx}%
\setcounter{MaxMatrixCols}{30}
%TCIDATA{OutputFilter=latex2.dll}
%TCIDATA{Version=5.50.0.2960}
%TCIDATA{CSTFile=revtex4.cst}
%TCIDATA{Created=Tuesday, November 25, 2014 10:35:54}
%TCIDATA{LastRevised=Friday, August 07, 2015 15:31:26}
%TCIDATA{<META NAME="GraphicsSave" CONTENT="32">}
%TCIDATA{<META NAME="SaveForMode" CONTENT="1">}
%TCIDATA{BibliographyScheme=Manual}
%TCIDATA{<META NAME="DocumentShell" CONTENT="Articles\SW\REVTeX 4">}
%BeginMSIPreambleData
\providecommand{\U}[1]{\protect\rule{.1in}{.1in}}
%EndMSIPreambleData
\newtheorem{theorem}{Theorem}
\newtheorem{acknowledgement}[theorem]{Acknowledgement}
\newtheorem{algorithm}[theorem]{Algorithm}
\newtheorem{axiom}[theorem]{Axiom}
\newtheorem{claim}[theorem]{Claim}
\newtheorem{conclusion}[theorem]{Conclusion}
\newtheorem{condition}[theorem]{Condition}
\newtheorem{conjecture}[theorem]{Conjecture}
\newtheorem{corollary}[theorem]{Corollary}
\newtheorem{criterion}[theorem]{Criterion}
\newtheorem{definition}[theorem]{Definition}
\newtheorem{example}[theorem]{Example}
\newtheorem{exercise}[theorem]{Exercise}
\newtheorem{lemma}[theorem]{Lemma}
\newtheorem{notation}[theorem]{Notation}
\newtheorem{problem}[theorem]{Problem}
\newtheorem{proposition}[theorem]{Proposition}
\newtheorem{remark}[theorem]{Remark}
\newtheorem{solution}[theorem]{Solution}
\newtheorem{summary}[theorem]{Summary}
\newenvironment{proof}[1][Proof]{\noindent\textbf{#1.} }{\ \rule{0.5em}{0.5em}}
\begin{document}
\preprint{HEP/123-qed}
\title[Short title for running header]{REV\TeX4}
\author{First Author}
\affiliation{My Institution}
\author{Second Author}
\affiliation{My Institution}
\author{Third Author}
\affiliation{Other Institution}
\keywords{one two three}
\pacs{PACS number}

\begin{abstract}
Shell document for REV\TeX{} 4.

\end{abstract}
\volumeyear{year}
\volumenumber{number}
\issuenumber{number}
\eid{identifier}
\date[Date text]{date}
\received[Received text]{date}

\revised[Revised text]{date}

\accepted[Accepted text]{date}

\published[Published text]{date}

\startpage{101}
\endpage{102}
\maketitle
\tableofcontents


\subsection{\bigskip$\left(  \text{Real}\right)  \dim\left\{  V_{2}\right\}
=2\protect\cong\mathbb{R}^{2};\left(  \text{Complex}\right)  \dim\left\{
V_{2}\right\}  =1$}

\paragraph{Basic Properties}

A complex unit in $V_{2}$ is given by
\begin{align}
J_{e}  &  =\left\vert e_{2}\right)  \left(  e_{1}\right\vert -\left\vert
e_{1}\right)  \left(  e_{2}\right\vert \\
R_{e}\left(  J_{e}\right)   &  =\left(
\begin{array}
[c]{cc}%
0 & -1\\
1 & 0
\end{array}
\right)
\end{align}
the associated complex conjugation
\begin{align}
\sigma_{e}  &  :\mathbb{R}^{2}\rightarrow\mathbb{R}^{2}\\
\sigma_{e}  &  =\left\vert e_{1}\right)  \left(  e_{1}\right\vert -\left\vert
e_{2}\right)  \left(  e_{2}\right\vert =\left\vert e_{1}\right)  \left(
e_{1}\right\vert -J_{e}\left\vert e_{1}\right)  \left(  e_{1}\right\vert
J_{e}^{t}\\
R_{e}\left(  \sigma_{e}\right)   &  =\left(
\begin{array}
[c]{cc}%
1 & 0\\
0 & -1
\end{array}
\right)
\end{align}
and the definition of real
\begin{align}
\left\vert a_{r}\right)   &  =\frac{1}{2}\left(  \left\vert a\right)
+\sigma_{e}\left(  \left\vert a\right)  \right)  \right)  =\frac{1}%
{2}\left\vert a\right)  +\frac{1}{2}\left(  \left\vert e_{1}\right)  \left(
e_{1}\right\vert -\left\vert e_{2}\right)  \left(  e_{2}\right\vert \right)
\left\vert a\right)  =\left\vert e_{1}\right)  a_{1}\\
&  =\frac{1}{2}\left\vert a\right)  +\frac{1}{2}\left(  \left\vert
e_{1}\right)  \left(  e_{1}\right\vert -J_{e}\left\vert e_{1}\right)  \left(
e_{1}\right\vert J_{e}^{t}\right)  \left\vert a\right) \\
R_{e}\left(  \left\vert a_{r}\right)  \right)   &  =\frac{1}{2}\left(  \left(
%
\begin{array}
[c]{c}%
a_{1}\\
a_{2}%
\end{array}
\right)  +\left(
\begin{array}
[c]{c}%
a_{1}\\
-a_{2}%
\end{array}
\right)  \right)  =\left(
\begin{array}
[c]{c}%
a_{1}\\
0
\end{array}
\right)
\end{align}
complex parts%
\begin{align}
\left\vert a_{i}\right)   &  =\frac{1}{2}J_{e}\left(  \left\vert a\right)
-\sigma_{e}\left(  \left\vert a\right)  \right)  \right)  =J_{e}\left(
\left\vert e_{1}\right)  \right)  a_{2}\\
R_{e}\left(  \left\vert a_{i}\right)  \right)   &  =\frac{1}{2}\left(
\begin{array}
[c]{cc}%
0 & -1\\
1 & 0
\end{array}
\right)  \left(  \left(
\begin{array}
[c]{c}%
a_{1}\\
a_{2}%
\end{array}
\right)  +\left(
\begin{array}
[c]{c}%
a_{1}\\
-a_{2}%
\end{array}
\right)  \right) \\
&  =\frac{1}{2}\left(  \left(
\begin{array}
[c]{c}%
-a_{2}\\
a_{1}%
\end{array}
\right)  +\left(
\begin{array}
[c]{c}%
a_{2}\\
a_{1}%
\end{array}
\right)  \right)  =\left(
\begin{array}
[c]{c}%
0\\
a_{2}%
\end{array}
\right)
\end{align}
and real vectors
\begin{align}
\sigma_{e}\left(  a\right)   &  =\left\{  \left\vert e_{1}\right)  \left(
e_{1}\right\vert -\left\vert e_{2}\right)  \left(  e_{2}\right\vert \right\}
a=\left\{  \left\vert e_{1}\right)  \left(  e_{1}\right\vert -J_{e}\left\vert
e_{1}\right)  \left(  e_{1}\right\vert J_{e}^{t}\right\}  a=a\\
&  =\left\vert e_{1}\right)  \left(  \left.  e_{1}\right\vert a\right)
-\left\vert e_{2}\right)  \left(  \left.  e_{2}\right\vert a\right)
=\left\vert e_{1}\right)  \left(  \left.  e_{1}\right\vert a\right)
+\left\vert e_{2}\right)  \left(  \left.  e_{2}\right\vert a\right) \\
&  \Rightarrow a=\left\vert e_{1}\right)  \left(  \left.  e_{1}\right\vert
a\right)  =\left\vert e_{1}\right)  a_{1}.
\end{align}
A representation of $\mathbb{C}$ in $\mathcal{B}\left(  V_{2}\right)  $ is
given by
\begin{equation}
K_{e}\left(  a\right)  =K_{e}\left(  R_{e}\left(  \left\vert a\right)
\right)  \right)  =\left(  a_{1}I_{2}+a_{2}J_{e}\right)  ;a_{1},a_{2}%
\in\mathbb{R},
\end{equation}
thus
\begin{align}
\mathfrak{C}_{e}\left(  \left\vert a\right)  \right)   &  =\mathfrak{C}%
_{e}\left(  \left(  \left\vert e_{1}\right)  \left(  e_{1}\right\vert
+\left\vert e_{2}\right)  \left(  e_{2}\right\vert \right)  \left\vert
a\right)  \right)  =\mathfrak{C}_{e}\left(  \left(  \left\vert e_{1}\right)
\left(  e_{1}\right\vert +J_{e}\left\vert e_{1}\right)  \left(  e_{1}%
\right\vert J_{e}^{t}\right)  \left\vert a\right)  \right) \\
&  =\left(  \left.  e_{1}\right\vert a\right)  \left\vert e_{1}\right)
+\left(  \left.  e_{2}\right\vert a\right)  J_{e}\left\vert e_{1}\right) \\
&  =\left\{  \left(  \left.  e_{1}\right\vert a\right)  I_{2}+J_{e}\left(
\left.  e_{2}\right\vert a\right)  \right\}  \left\vert e_{1}\right) \\
&  =K_{e}\left(  a\right)  \left\vert e_{1}\right)  .
\end{align}


\paragraph{Complexification Maps}

The inverse vector complexification map is a real linear map
\begin{equation}
\mathfrak{C}_{e}^{-1}:V_{2\mathbb{C}}\rightarrow V_{2}%
\end{equation}%
\begin{align}
\mathfrak{C}_{e}^{-1}\left(  \left\vert e_{1}\right\rangle \right)   &
=\left\vert e_{1}\right) \\
\mathfrak{C}_{e}^{-1}\left(  i\left\vert e_{1}\right\rangle \right)   &
=J_{e}\left\vert e_{1}\right)  =\left\vert e_{2}\right)
\end{align}
and the inverse operator complexification map
\begin{equation}
\widehat{\mathfrak{C}_{e}}^{-1}\left(  X\right)  =\mathfrak{C}_{e}%
^{-1}X\mathfrak{C}_{e}%
\end{equation}
\begin{align}
\widehat{\mathfrak{C}_{e}}^{-1}\left(  I_{V_{2\mathbb{C}}}\right)   &
=I_{V_{2}}=\left\vert e_{1}\right)  \left(  e_{1}\right\vert +\left\vert
e_{2}\right)  \left(  e_{2}\right\vert \\
\widehat{\mathfrak{C}_{e}}^{-1}\left(  iI_{V_{2\mathbb{C}}}\right)   &
=\mathfrak{C}_{e}^{-1}\left(  i\left\vert e_{1}\right\rangle \left\langle
e_{1}\right\vert \right)  =J_{e}=\left\vert e_{2}\right)  \left(
e_{1}\right\vert -\left\vert e_{1}\right)  \left(  e_{2}\right\vert
\end{align}
gives%
\begin{align}
\widehat{\mathfrak{C}_{e}}^{-1}\left(  zI_{V_{2\mathbb{C}}}\right)   &
=\widehat{\mathfrak{C}_{e}}^{-1}\left(  a_{1r}I_{V_{2\mathbb{C}}}%
+a_{1i}iI_{V_{2\mathbb{C}}}\right)  =a_{1r}I_{V_{2}}+a_{1i}J_{e}\\
R_{e}\left(  a_{1r}I_{V_{2}}+a_{1i}J_{e}\right)   &  =\left(
\begin{array}
[c]{cc}%
a_{1r} & -a_{1i}\\
a_{1i} & a_{1r}%
\end{array}
\right)  .
\end{align}
The opertor complexification
\begin{equation}
\widehat{\mathfrak{C}_{e}}\left(  X\right)  =\mathfrak{C}_{e}X\mathfrak{C}%
_{e}^{-1}%
\end{equation}
is given by
\begin{align}
\widehat{\mathfrak{C}_{e}}\left(  I_{V_{2}}\right)   &  =\mathfrak{C}%
_{e}\left(  \left\vert e_{1}\right)  \left(  e_{1}\right\vert +\left\vert
e_{2}\right)  \left(  e_{2}\right\vert \right)  =I_{V_{2\mathbb{C}}}\\
\widehat{\mathfrak{C}_{e}}\left(  J_{e}\right)   &  =\mathfrak{C}_{e}\left(
\left\vert e_{2}\right)  \left(  e_{1}\right\vert -\left\vert e_{1}\right)
\left(  e_{2}\right\vert \right)  =iI_{V_{2\mathbb{C}}}=i\left\vert
e_{1}\right\rangle \left\langle e_{1}\right\vert .
\end{align}


The Hilbert space $V_{2\mathbb{C}}$ has an associated real bilinear form that
is independent of the basis of $V_{2\mathbb{C}}$
\begin{align}
\left(  \left.  \cdot\right\vert \cdot\right)  _{V_{2\mathbb{C}}}  &
:V_{2\mathbb{C}}\times V_{2\mathbb{C}}\rightarrow\mathbb{R}\\
\left(  \left.  z_{1}\right\vert z_{2}\right)  _{V_{2\mathbb{C}}}  &
=\operatorname{Re}\left\langle \left.  z_{1}\right\vert z_{2}\right\rangle
\end{align}
with the properties
\begin{align}
\left(  \left.  e_{1}\right\vert ie_{1}\right)  _{V_{2\mathbb{C}}}  &  =0\\
\left(  \left.  e_{1}\right\vert e_{1}\right)  _{V_{2\mathbb{C}}}  &  =1\\
\left(  \left.  ie_{1}\right\vert ie_{1}\right)  _{V_{2\mathbb{C}}}  &  =1
\end{align}
and is non degenerate i.e.
\begin{equation}
\left(  \left.  z\right\vert z\right)  _{V_{2\mathbb{C}}}=0\Leftrightarrow
\left\vert z\right\rangle =0.
\end{equation}
The form $\left(  \left.  \cdot\right\vert \cdot\right)  _{V_{2\mathbb{C}}}$
induces a real linear functional, $\left(  z\right\vert _{V_{2\mathbb{C}}},$
on $V_{2\mathbb{C}}$ defined by
\begin{align}
\left(  z\right\vert _{V_{2\mathbb{C}}}  &  :V_{2\mathbb{C}}\rightarrow
\mathbb{R}\\
\left(  z\right\vert _{V_{2\mathbb{C}}}\left(  \left\vert \omega\right\rangle
\right)   &  =\left(  \left.  z\right\vert \left\vert \omega\right\rangle
\right)  _{V_{2\mathbb{C}}}=\operatorname{Re}\left\langle \left.  z\right\vert
\omega\right\rangle
\end{align}
i.e.
\begin{align}
\left(  z\right\vert _{V_{2\mathbb{C}}}\left(  \alpha_{1}\left\vert \omega
_{1}\right\rangle +\alpha_{2}\left\vert \omega_{2}\right\rangle \right)   &
=\alpha_{1}\left(  z\right\vert _{V_{2\mathbb{C}}}\left(  \left\vert
\omega_{1}\right\rangle \right)  +\alpha_{2}\left(  z\right\vert
_{V_{2\mathbb{C}}}\left(  \left\vert \omega_{2}\right\rangle \right)
;\alpha_{1},\alpha_{2}\in\mathbb{R}\\
\left(  z\right\vert _{V_{2\mathbb{C}}}\left(  \alpha_{1}\left\vert \omega
_{1}\right\rangle +\alpha_{2}\left\vert \omega_{2}\right\rangle \right)   &
\neq\alpha_{1}\left(  z\right\vert _{V_{2\mathbb{C}}}\left(  \left\vert
\omega_{1}\right\rangle \right)  +\alpha_{2}\left(  z\right\vert
_{V_{2\mathbb{C}}}\left(  \left\vert \omega_{2}\right\rangle \right)
;\alpha_{1},\alpha_{2}\in\mathbb{C}.
\end{align}


\begin{theorem}%
\begin{equation}
T:\left\vert z\right\rangle \rightarrow\left(  z\right\vert _{V_{2\mathbb{C}}}%
\end{equation}
is injective i.e.
\begin{equation}
T\left(  \left\vert z_{1}\right\rangle -\left\vert z_{2}\right\rangle \right)
=0\Leftrightarrow\left\vert z_{1}\right\rangle =\left\vert z_{2}\right\rangle
\end{equation}


\begin{proof}%
\begin{equation}
\left(  \left.  z\right\vert z\right)  _{V_{2\mathbb{C}}}=\operatorname{Re}%
\left\langle \left.  z\right\vert z\right\rangle =\left\langle \left.
z\right\vert z\right\rangle
\end{equation}
thus if
\begin{equation}
\left(  \left.  \left(  z_{1}-z_{2}\right)  \right\vert z\right)
=0\forall\left\vert z_{2}\right\rangle \in V_{2\mathbb{C}}%
\end{equation}
then in particular
\begin{equation}
\left(  \left.  \left(  z_{1}-z_{2}\right)  \right\vert \left(  z_{1}%
-z_{2}\right)  \right)  =\left\langle \left.  \left(  z_{1}-z_{2}\right)
\right\vert \left(  z_{1}-z_{2}\right)  \right\rangle =0\Rightarrow
z_{1}=z_{2}%
\end{equation}
and if
\begin{equation}
z_{1}=z_{2}%
\end{equation}
then
\begin{equation}
\left(  \left.  \left(  z_{1}-z_{2}\right)  \right\vert z\right)
=0\forall\left\vert z_{2}\right\rangle \in V_{2\mathbb{C}}.
\end{equation}

\end{proof}
\end{theorem}

\paragraph{Dual Space Representation}

A complexification map and its inverse can then be expressed as the real
linear maps
\begin{equation}
\mathfrak{C}_{e}=\left\vert e_{1}\right\rangle \left(  e_{1}\right\vert
+\left\vert ie_{1}\right\rangle \left(  e_{2}\right\vert =\left(  \left\vert
e_{1}\right\rangle ,\left\vert ie_{1}\right\rangle \right)  \left(
\begin{array}
[c]{cc}%
1 & 0\\
0 & 1
\end{array}
\right)  \left(
\begin{array}
[c]{c}%
\left(  e_{1}\right\vert \\
\left(  e_{2}\right\vert
\end{array}
\right)
\end{equation}
and%
\begin{equation}
\mathfrak{C}_{e}^{-1}=\left\vert e_{1}\right)  \left(  e_{1}\right\vert
_{V_{2\mathbb{C}}}+\left\vert e_{2}\right)  \left(  ie_{1}\right\vert
_{V_{2\mathbb{C}}}=\left\vert e_{1}\right)  \left(  e_{1}\right\vert
_{V_{2\mathbb{C}}}+J_{e}\left\vert e_{1}\right)  \left(  ie_{1}\right\vert
_{V_{2\mathbb{C}}}=\left(  \left\vert e_{1}\right)  ,\left\vert e_{2}\right)
\right)  \left(
\begin{array}
[c]{cc}%
1 & 0\\
0 & 1
\end{array}
\right)  \left(
\begin{array}
[c]{c}%
\left(  e_{1}\right\vert _{V_{2\mathbb{C}}}\\
\left(  ie_{1}\right\vert _{V_{2\mathbb{C}}}%
\end{array}
\right)  ,
\end{equation}
which give with respect to the real structure of $V_{2\mathbb{C}}$
\begin{align}
\mathfrak{C}_{e}\left\vert e_{1},e_{2}\right)   &  =\left\vert \mathfrak{C}%
_{e}\left(  e_{1}\right)  ,\mathfrak{C}_{e}\left(  e_{2}\right)  \right\rangle
=\left\vert e_{1},\mathfrak{i}e_{1}\right\rangle \left(
\begin{array}
[c]{cc}%
1 & 0\\
0 & 1
\end{array}
\right)  =\left\vert e_{1},\mathfrak{i}e_{1}\right\rangle R_{e}\left(
\mathfrak{C}_{e}\right) \\
\mathfrak{C}_{e}^{-1}\left\vert e_{1},\mathfrak{i}e_{1}\right\rangle  &
=\left\vert \mathfrak{C}_{e}^{-1}\left(  e_{1}\right)  ,\mathfrak{C}_{e}%
^{-1}\left(  \mathfrak{i}e_{1}\right)  \right)  =\left\vert e_{1}%
,e_{2}\right)  \left(
\begin{array}
[c]{cc}%
1 & 0\\
0 & 1
\end{array}
\right)  =\left\vert e_{1},e_{2}\right)  R_{e}\left(  \mathfrak{C}_{e}%
^{-1}\right)  .
\end{align}
This leads to
\begin{align}
\mathfrak{C}_{e}\mathfrak{C}_{e}^{-1}  &  =\left(  \left\vert e_{1}%
\right\rangle \left(  e_{1}\right\vert +\left\vert ie_{1}\right\rangle \left(
e_{2}\right\vert \right)  \left(  \left\vert e_{1}\right)  \left(
e_{1}\right\vert _{V_{2\mathbb{C}}}+\left\vert e_{2}\right)  \left(
ie_{1}\right\vert _{V_{2\mathbb{C}}}\right) \\
&  =\left\vert e_{1}\right\rangle \left(  e_{1}\right\vert _{V_{2\mathbb{C}}%
}+0+0+\left\vert ie_{1}\right\rangle \left(  ie_{1}\right\vert
_{V_{2\mathbb{C}}}=\left\vert e_{1}\right\rangle \left(  e_{1}\right\vert
_{V_{2\mathbb{C}}}+\left\vert ie_{1}\right\rangle \left(  ie_{1}\right\vert
_{V_{2\mathbb{C}}}=I_{V_{2\mathbb{C}}}%
\end{align}
and%
\begin{equation}
\mathfrak{C}_{e}^{-1}\mathfrak{C}_{e}=\left(  \left\vert e_{1}\right)  \left(
e_{1}\right\vert _{V_{2\mathbb{C}}}+\left\vert e_{2}\right)  \left(
ie_{1}\right\vert _{V_{2\mathbb{C}}}\right)  \left(  \left\vert e_{1}%
\right\rangle \left(  e_{1}\right\vert +\left\vert ie_{1}\right\rangle \left(
e_{2}\right\vert \right)  =\left\vert e_{1}\right)  \left(  e_{1}\right\vert
+\left\vert e_{2}\right)  \left(  e_{2}\right\vert =I_{V_{2}}.
\end{equation}


The complex unit, $J_{e_{\mathbb{C}}},$ $\left(  \text{which is real
linear}\right)  $ associated with a basis $\left\{  \left\vert \mathbf{e}%
\right\rangle \right\}  $ can be represented in $V_{2\mathbb{C}}$ by
\begin{align}
J_{e_{\mathbb{C}}}  &  =\left\vert ie_{1}\right\rangle \left(  e_{1}%
\right\vert _{V_{2\mathbb{C}}}-\left\vert e_{1}\right\rangle \left(
ie_{1}\right\vert _{V_{2\mathbb{C}}}\\
J_{e_{\mathbb{C}}}\left\vert z_{1}\right\rangle  &  =\left(  a_{11}%
+ia_{12}\right)  \left\vert e_{1}\right\rangle =i\left(  a_{11}+ia_{12}%
\right)  \left\vert e_{1}\right\rangle =i\left\vert z_{1}\right\rangle .
\end{align}
The complex conjugation map , $\sigma_{e_{\mathbb{C}}},$ $\left(  \text{which
is real linear}\right)  $ associated with this basis can be expressed as
\[
\sigma_{e_{\mathbb{C}}}=\left\vert e_{1}\right\rangle \left(  e_{1}\right\vert
_{V_{2\mathbb{C}}}-\left\vert ie_{1}\right\rangle \left(  ie_{1}\right\vert
_{V_{2\mathbb{C}}}%
\]
and thus the definition of $\mathbf{e}$-real vectors becomes
\begin{equation}
\sigma_{e_{\mathbb{C}}}\left(  \left\vert z\right\rangle \right)  =\left\vert
z\right\rangle .
\end{equation}


The operator%

\begin{equation}
\left\vert e_{1}\right\rangle \left(  e_{1}\right\vert _{V_{2\mathbb{C}}%
}+\left\vert ie_{1}\right\rangle \left(  ie_{1}\right\vert _{V_{2\mathbb{C}}}%
\end{equation}
resolves vectors $\left\vert z_{1}\right\rangle \in V_{2\mathbb{C}}$ into
their real and complex parts wrt the basis $\left\{  \left\vert e_{1}%
\right\rangle \right\}  $ i.e.
\begin{equation}
\left(  \left\vert e_{1}\right\rangle \left(  e_{1}\right\vert
_{V_{2\mathbb{C}}}+\left\vert ie_{1}\right\rangle \left(  ie_{1}\right\vert
_{V_{2\mathbb{C}}}\right)  \left\vert z_{1}\right\rangle =\left\vert
e_{1}\right\rangle \left(  \left.  e_{1}\right\vert z_{1}\right\rangle
+\left\vert ie_{1}\right\rangle \left(  \left.  ie_{1}\right\vert
z_{1}\right\rangle =\left\vert e_{1}\right\rangle a_{11}+\left\vert
ie_{1}\right\rangle a_{12}=\left\vert e_{1},ie_{1}\right\rangle R_{e}\left(
a_{1}\right)  \equiv\left\vert e_{1},ie_{1}\right\rangle R_{ee}\left(
z_{1}\right)  .
\end{equation}


\subsection{Detailed Example of Complexification}

\paragraph{$\left(  \text{Real}\right)  \dim\left\{  V_{2}\right\}
=2\protect\cong\mathbb{R}^{2}$}

In order to understand a complexification, $V_{2s\mathbb{C}},$ of a real even
dimensional vector space $V_{2s},$ we first consider
\begin{equation}
\left(  \text{real}\right)  \dim\left\{  V_{2}\right\}  =2;\left(
\text{complex}\right)  \dim\left\{  V_{2\mathbb{C}}\right\}  =1,
\end{equation}
where
\begin{align}
\left\vert e_{1}\right)   &  =\left\vert \mathbf{e}\right)  \left(
\begin{array}
[c]{c}%
1\\
0
\end{array}
\right) \\
\left\vert e_{2}\right)   &  =\left\vert \mathbf{e}\right)  \left(
\begin{array}
[c]{c}%
0\\
1
\end{array}
\right)
\end{align}
and the complexification, $J_{e},$ is given by%
\begin{equation}
J_{e}\left\vert e_{1}\right)  =\left\vert \mathbf{e}\right)  \left(
\begin{array}
[c]{cc}%
0 & -1\\
1 & 0
\end{array}
\right)  \left(
\begin{array}
[c]{c}%
1\\
0
\end{array}
\right)  =\left\vert \mathbf{e}\right)  \left(
\begin{array}
[c]{c}%
0\\
1
\end{array}
\right)  =\left\vert e_{2}\right)
\end{equation}
%

\begin{equation}
J_{e}\left\vert e_{2}\right)  =\left\vert \mathbf{e}\right)  \left(
\begin{array}
[c]{cc}%
0 & -1\\
1 & 0
\end{array}
\right)  \left(
\begin{array}
[c]{c}%
0\\
1
\end{array}
\right)  =\left\vert \mathbf{e}\right)  \left(
\begin{array}
[c]{c}%
-1\\
0
\end{array}
\right)  =-\left\vert e_{1}\right)  .
\end{equation}
Hence
\begin{equation}
J_{e}=\left\vert e_{2}\right)  \left(  e_{1}\right\vert -\left\vert
e_{1}\right)  \left(  e_{2}\right\vert ,
\end{equation}
so corresponds to a real antisymmetric bilinear form, $g,$%
\begin{align}
g\left(  a,b\right)   &  =\left(  \left.  g\right\vert a\wedge b\right)
=\left(  \left.  a\right\vert J_{e}b\right) \\
\left\vert g\right)   &  =\left\vert e_{1}e_{2}\right)
\end{align}
on $V_{2},$ while
\begin{equation}
I_{2}=\left\vert e_{1}\right)  \left(  e_{1}\right\vert +\left\vert
e_{2}\right)  \left(  e_{2}\right\vert .
\end{equation}


This real space can be thought of a complex space of half the dimension
through a complexification map
\begin{equation}
\mathfrak{C}_{e}=\left\vert e_{1}\right\rangle \left(  e_{1}\right\vert
+\left\vert ie_{2}\right\rangle \left(  e_{2}\right\vert
\end{equation}
giving%
\begin{align}
\mathfrak{C}_{e}\left(  \left\vert a\right)  \right)   &  =\mathfrak{C}%
_{e}\left(  \left\vert e_{1}\right)  a_{11}+\left\vert e_{2}\right)
a_{12}\right)  =\left\vert e_{1}\right\rangle a_{11}+\left\vert ie_{1}%
\right\rangle a_{12}\\
&  =\left\vert e_{1}\right\rangle a_{11}+\left\vert e_{1}\right\rangle
ia_{12}=\left\vert e_{1}\right\rangle z_{1}=\left\vert e_{1}\right\rangle
R_{e}\left(  z\right)
\end{align}


This complexification map
\begin{equation}
\mathfrak{C}_{e}:V_{2}\rightarrow V_{2e\mathbb{C}},
\end{equation}
can also be expressed through the operator complexification, where%
\begin{align}
\mathfrak{C}_{e}\left(  \left\vert e_{1}\right)  \right)   &  =\mathfrak{C}%
_{e}\left(  \left(  1I_{2}+0J_{e}\right)  \left\vert e_{1}\right)  \right)
=\widehat{\mathfrak{C}}_{e}\left(  1I_{2}+0J_{e}\right)  \mathfrak{C}%
_{e}\left(  \left\vert e_{1}\right)  \right)  =\left\vert e_{1}\right\rangle
\\
\mathfrak{C}_{e}\left(  \left\vert e_{2}\right)  \right)   &  =\mathfrak{C}%
_{e}\left(  \left(  0I_{2}+1J_{e}\right)  \left\vert e_{1}\right)  \right)
=\widehat{\mathfrak{C}}_{e}\left(  0I_{2}+1J_{e}\right)  \mathfrak{C}%
_{e}\left(  \left\vert e_{1}\right)  \right)  =i\left\vert e_{1}\right\rangle
\end{align}
giving%
\begin{align}
\mathfrak{C}_{e}\left(  \left\vert a\right)  \right)   &  =\mathfrak{C}%
_{e}\left(  \left\vert e_{1}\right)  a_{11}+\left\vert e_{2}\right)
a_{12}\right)  =\\
\mathfrak{C}_{e}\left(  \left(  a_{11}I_{2}+a_{12}J_{e}\right)  \left\vert
e_{1}\right)  \right)   &  =\left(  a_{11}+a_{12}i\right)  \left\vert
e_{1}\right\rangle =z_{1}\left\vert e_{1}\right\rangle =\left\vert
z\right\rangle ,
\end{align}
i.e.
\begin{equation}
\mathfrak{C}_{e}\left(  \left\vert a\right)  \right)  =\mathfrak{C}_{e}\left(
K_{e}\left(  z_{1}\right)  \left\vert e_{1}\right)  \right)  =\left\vert
z\right\rangle =z_{1}\left\vert e_{1}\right\rangle .
\end{equation}
A complexifiaction and its inverse are given by
\begin{equation}
\mathfrak{C}_{e}=\left\vert e_{1}\right\rangle \left(  e_{1}\right\vert
+\left\vert ie_{1}\right\rangle \left(  e_{2}\right\vert =\left(  \left\vert
e_{1}\right\rangle ,\left\vert ie_{1}\right\rangle \right)  \left(
\begin{array}
[c]{cc}%
1 & 0\\
0 & 1
\end{array}
\right)  \left(
\begin{array}
[c]{c}%
\left(  e_{1}\right\vert \\
\left(  e_{2}\right\vert
\end{array}
\right)
\end{equation}
and%
\begin{equation}
\mathfrak{C}_{e}^{-1}=\left\vert e_{1}\right)  \left(  e_{1}\right\vert
_{V_{2\mathbb{C}}}+\left\vert e_{2}\right)  \left(  ie_{1}\right\vert
_{V_{2\mathbb{C}}}=\left\vert e_{1}\right)  \left(  e_{1}\right\vert
_{V_{2\mathbb{C}}}+J_{e}\left\vert e_{1}\right)  \left(  ie_{1}\right\vert
_{V_{2\mathbb{C}}}=\left(  \left\vert e_{1}\right)  ,\left\vert e_{2}\right)
\right)  \left(
\begin{array}
[c]{cc}%
1 & 0\\
0 & 1
\end{array}
\right)  \left(
\begin{array}
[c]{c}%
\left(  e_{1}\right\vert _{V_{2\mathbb{C}}}\\
\left(  ie_{1}\right\vert _{V_{2\mathbb{C}}}%
\end{array}
\right)  ,
\end{equation}
Thus
\begin{align}
\mathfrak{C}_{e}  &  =\left\vert e_{1}\right\rangle \left(  e_{1}\right\vert
+i\left\vert e_{1}\right\rangle \left(  J_{e}e_{1}\right\vert \\
\mathfrak{C}_{e}\left(  \left\vert a\right)  \right)   &  =\left(  \left\vert
e_{1}\right\rangle \left(  e_{1}\right\vert +i\left\vert e_{1}\right\rangle
\left(  J_{e}e_{1}\right\vert \right)  \left\vert a\right)  =\left\vert
e_{1}\right\rangle a_{11}+i\left\vert e_{1}\right\rangle a_{12}=\left(
a_{11}+ia_{12}\right)  \left\vert e_{1}\right\rangle =z_{1}\left\vert
e_{1}\right\rangle =\left\vert z\right\rangle \\
\mathfrak{C}_{e}\left(  \left\vert e_{1}\right)  \right)   &  =\left\vert
e_{1}\right\rangle \\
\mathfrak{C}_{e}\left(  \left\vert e_{2}\right)  \right)   &  =\mathfrak{C}%
_{e}\left(  \left\vert J_{e}e_{1}\right)  \right)  =i\left\vert e_{1}%
\right\rangle
\end{align}
The inverse map
\begin{align}
\mathfrak{C}_{e}^{-1}  &  :V_{2e\mathbb{C}}\rightarrow V_{2},\\
\mathfrak{C}_{e}^{-1}  &  =\left\vert e_{1}\right)  \left(  e_{1}\right\vert
_{V_{2\mathbb{C}}}+\left\vert e_{2}\right)  \left(  ie_{1}\right\vert
_{V_{2\mathbb{C}}}=\left\vert e_{1}\right)  \left(  e_{1}\right\vert
_{V_{2\mathbb{C}}}+J_{e}\left\vert e_{1}\right)  \left(  ie_{1}\right\vert
_{V_{2\mathbb{C}}}\\
\mathfrak{C}_{e}^{-1}\left(  \left\vert z\right\rangle \right)   &  =\left(
\left\vert e_{1}\right)  \left(  e_{1}\right\vert _{V_{2\mathbb{C}}%
}+\left\vert e_{2}\right)  \left(  ie_{1}\right\vert _{V_{2\mathbb{C}}%
}\right)  \left\vert z\right\rangle =\left\vert e_{1}\right)  a_{11}%
+\left\vert e_{2}\right)  a_{12}.
\end{align}


The products of a complexification map and its inverse give the expected
results i.e.
\begin{align}
\mathfrak{C}_{e}\mathfrak{C}_{e}^{-1}  &  =\left(  \left\vert e_{1}%
\right\rangle \left(  e_{1}\right\vert +i\left\vert e_{1}\right\rangle \left(
e_{2}\right\vert \right)  \left(  \left\vert e_{1}\right)  \left(
e_{1}\right\vert _{V_{2\mathbb{C}}}+\left\vert e_{2}\right)  \left(
ie_{1}\right\vert _{V_{2\mathbb{C}}}\right)  =I_{V_{2\mathbb{C}}}=\left\vert
e_{1}\right\rangle \left\langle e_{1}\right\vert =1\\
\mathfrak{C}_{e}^{-1}\mathfrak{C}_{e}  &  =\left(  \left\vert e_{1}\right)
\left(  e_{1}\right\vert _{V_{2\mathbb{C}}}+\left\vert e_{2}\right)  \left(
ie_{1}\right\vert _{V_{2\mathbb{C}}}\right)  \left(  \left\vert e_{1}%
\right\rangle \left(  e_{1}\right\vert +i\left\vert e_{1}\right\rangle \left(
e_{2}\right\vert \right)  =\left\vert e_{1}\right)  \left(  e_{1}\right\vert
+\left\vert e_{2}\right)  \left(  e_{2}\right\vert =I_{V_{2}}.
\end{align}


\subsubsection{Linear Transformation}

The relationship between complex linear transformation
\begin{equation}
T_{C}:V_{2e\mathbb{C}}\rightarrow V_{2e\mathbb{C}}%
\end{equation}
and real linear transformation
\begin{equation}
T_{R}:V_{2}\rightarrow V_{2}%
\end{equation}
is given by%
\begin{equation}
T_{C}=\mathfrak{C}_{e}T_{R}\mathfrak{C}_{e}^{-1}=\widehat{\mathfrak{C}}%
_{e}\left(  T_{R}\right)
\end{equation}


The action of $T_{C}$ is given by
\begin{equation}
\left\vert w\right\rangle =T_{C}\left\vert z\right\rangle =T_{C}\left\vert
e_{1}\right\rangle R_{e}\left(  z\right)  =\left\vert e_{1}\right\rangle
R_{e}\left(  T_{C}\right)  R_{e}\left(  z\right)  =\left\vert e_{1}%
\right\rangle tz,
\end{equation}
which can also be expressed as
\begin{equation}
\left\vert w\right\rangle =\mathfrak{C}_{e}\left(  b\right)  =\mathfrak{C}%
_{e}T_{R}\mathfrak{C}_{e}^{-1}\mathfrak{C}_{e}\left(  a\right)  =\mathfrak{C}%
_{e}\left(  T_{R}a\right)  .
\end{equation}
In more detail%
\begin{align}
\left\vert w\right\rangle  &  =T_{C}\left\vert z_{1}\right\rangle
=T_{C}\left\vert e_{1}\right\rangle R_{e}\left(  z_{1}\right)  =\left\vert
e_{1}\right\rangle R_{e}\left(  T_{C}\right)  R_{e}\left(  z_{1}\right) \\
R_{e}\left(  T_{C}\right)  R_{e}\left(  z_{1}\right)   &  =\left(
\alpha+i\beta\right)  \left(  x+iy\right)  =\left(  \alpha x-\beta y\right)
+i\left(  \beta x+\alpha y\right)  ,
\end{align}
which can be expressed in terms of the over complete basis $\left\{
\left\vert e_{1}\right\rangle ,\left\vert ie_{1}\right\rangle \right\}  $ of
$V_{2\mathbb{C}}$ using
\begin{align}
T_{C}  &  =\left(
\begin{array}
[c]{cc}%
\left\vert e_{1}\right\rangle  & \left\vert ie_{1}\right\rangle
\end{array}
\right)  \left(
\begin{array}
[c]{cc}%
\alpha & -\beta\\
\beta & \alpha
\end{array}
\right)  \left(
\begin{array}
[c]{c}%
\left(  e_{1}\right\vert _{V_{2\mathbb{C}}}\\
\left(  ie_{1}\right\vert _{V_{2\mathbb{C}}}%
\end{array}
\right) \\
&  =\left(  \left\vert e_{1}\right\rangle \left(  e_{1}\right\vert
_{V_{2\mathbb{C}}}+\left\vert ie_{1}\right\rangle \left(  ie_{1}\right\vert
_{V_{2\mathbb{C}}}\right)  R_{ee}\left(  T_{C}\right)  \left(  \left\vert
e_{1}\right\rangle \left(  e_{1}\right\vert _{V_{2\mathbb{C}}}+\left\vert
ie_{1}\right\rangle \left(  ie_{1}\right\vert _{V_{2\mathbb{C}}}\right) \\
&  =\left\vert e_{1}\right\rangle \left(  e_{1}\right\vert _{V_{2\mathbb{C}}%
}T_{C}\left\vert e_{1}\right\rangle \left(  e_{1}\right\vert _{V_{2\mathbb{C}%
}}+\left\vert e_{1}\right\rangle \left(  e_{1}\right\vert _{V_{2\mathbb{C}}%
}T_{C}\left\vert ie_{1}\right\rangle \left(  ie_{1}\right\vert
_{V_{2\mathbb{C}}}+\left\vert ie_{1}\right\rangle \left(  ie_{1}\right\vert
_{V_{2\mathbb{C}}}T_{C}\left\vert e_{1}\right\rangle \left(  e_{1}\right\vert
_{V_{2\mathbb{C}}}+\left\vert ie_{1}\right\rangle \left(  ie_{1}\right\vert
_{V_{2\mathbb{C}}}T_{C}\left\vert ie_{1}\right\rangle \left(  ie_{1}%
\right\vert _{V_{2\mathbb{C}}}\\
&  =\alpha\left\vert e_{1}\right\rangle \left(  e_{1}\right\vert
_{V_{2\mathbb{C}}}-\beta\left\vert e_{1}\right\rangle \left(  ie_{1}%
\right\vert _{V_{2\mathbb{C}}}+\beta\left\vert ie_{1}\right\rangle \left(
e_{1}\right\vert _{V_{2\mathbb{C}}}+\alpha\left\vert ie_{1}\right\rangle
\left(  ie_{1}\right\vert _{V_{2\mathbb{C}}}%
\end{align}
giving
\begin{align}
T_{C}\left\vert z\right\rangle  &  =T_{C}\left\vert e_{1},ie_{1}\right\rangle
R_{ee}\left(  z_{1}\right)  =\left\vert e_{1},ie_{1}\right\rangle
R_{ee}\left(  T_{C}\right)  R_{ee}\left(  z_{1}\right)  =T_{C}\left\vert
e_{1},ie_{1}\right\rangle \left(
\begin{array}
[c]{c}%
a_{11}\\
a_{12}%
\end{array}
\right) \\
&  =\left\vert e_{1},ie_{1}\right\rangle \left(
\begin{array}
[c]{cc}%
\alpha & -\beta\\
\beta & \alpha
\end{array}
\right)  \left(
\begin{array}
[c]{c}%
a_{11}\\
a_{12}%
\end{array}
\right)  =\left\vert e_{1},ie_{1}\right\rangle \left(
\begin{array}
[c]{c}%
\alpha a_{11}-\beta a_{12}\\
\beta a_{11}+\alpha a_{12}%
\end{array}
\right)  =\left\vert e_{1},ie_{1}\right\rangle \left(
\begin{array}
[c]{c}%
b_{11}\\
b_{12}%
\end{array}
\right)  .
\end{align}


Applying the inverse complexification map gives%
\begin{align}
\widehat{\mathfrak{C}}_{e}^{-1}\left(  T_{C}\right)  \mathfrak{C}_{e}%
^{-1}\left(  \left\vert z_{1}\right\rangle \right)   &  =\widehat{\mathfrak{C}%
}_{e}^{-1}\left(  T_{C}\right)  \mathfrak{C}_{e}^{-1}\left(  \left\vert
z_{1}\right\rangle \right)  =\widehat{\mathfrak{C}}_{e}^{-1}\left(
T_{C}\right)  \left\vert \mathfrak{C}_{e}^{-1}\left(  e_{1}\right)
,\mathfrak{C}_{e}^{-1}\left(  ie_{1}\right)  \right\rangle \left(
\begin{array}
[c]{c}%
a_{11}\\
a_{12}%
\end{array}
\right) \\
&  =\widehat{\mathfrak{C}}_{e}^{-1}\left(  T_{C}\right)  \left\vert
e_{1},e_{2}\right)  \left(
\begin{array}
[c]{c}%
a_{11}\\
a_{12}%
\end{array}
\right)  =\left\vert e_{1},e_{2}\right)  \left(
\begin{array}
[c]{cc}%
\alpha & -\beta\\
\beta & \alpha
\end{array}
\right)  \left(
\begin{array}
[c]{c}%
a_{11}\\
a_{12}%
\end{array}
\right)  =\left\vert e_{1},e_{2}\right)  \left(
\begin{array}
[c]{c}%
\alpha a_{11}-\beta a_{12}\\
\beta a_{11}+\alpha a_{12}%
\end{array}
\right)  =\left\vert e_{1},e_{2}\right)  \left(
\begin{array}
[c]{c}%
b_{11}\\
b_{12}%
\end{array}
\right)  .
\end{align}
Hence
\begin{equation}
R_{e}\left(  \widehat{\mathfrak{C}}_{e}^{-1}\left(  t\right)  \right)
=\left(
\begin{array}
[c]{cc}%
\alpha & -\beta\\
\beta & \alpha
\end{array}
\right)  =R_{ee}\left(  t\right)  .
\end{equation}


The action of $T_{R}$ is given by%
\begin{equation}
\left\vert b\right)  =\mathfrak{C}_{e}^{-1}\left(  \left\vert w\right\rangle
\right)  =\widehat{\mathfrak{C}}_{e}^{-1}\left(  T_{C}\right)  \mathfrak{C}%
_{e}^{-1}\left\vert z\right\rangle =\mathfrak{C}_{e}^{-1}T_{C}\mathfrak{C}%
_{e}\mathfrak{C}_{e}^{-1}\left\vert z\right\rangle ,
\end{equation}
which can also be expressed as%
\begin{equation}
\left\vert b\right)  =T_{R}\left\vert a\right)  .
\end{equation}


The representations of $T_{C}$ and $T_{R}$ are related by
\begin{align}
R_{e}\left(  T_{R}\right)   &  =R_{e}\left(  \mathfrak{C}_{e}^{-1}%
T_{C}\mathfrak{C}_{e}\right)  =\left(
\begin{array}
[c]{cc}%
\alpha & -\beta\\
\beta & \alpha
\end{array}
\right)  =R_{e}\left(  \mathfrak{C}_{e}^{-1}\right)  R_{ee}\left(
T_{C}\right)  R_{e}\left(  \mathfrak{C}_{e}\right)  =R_{e}\left(
\mathfrak{C}_{e}^{-1}\right)  R_{ee}\left(  t\right)  R_{e}\left(
\mathfrak{C}_{e}\right) \\
R_{ee}\left(  T_{C}\right)   &  =R_{e}\left(  \mathfrak{C}_{e}\right)
R_{e}\left(  T_{R}\right)  R_{e}\left(  \mathfrak{C}_{e}^{-1}\right)
=R_{e}\left(  \mathfrak{C}_{e}\right)  \left(
\begin{array}
[c]{cc}%
\alpha & -\beta\\
\beta & \alpha
\end{array}
\right)  R_{e}\left(  \mathfrak{C}_{e}^{-1}\right)  ,
\end{align}
where
\begin{equation}
\mathfrak{C}_{e}=\left\vert e_{1}\right\rangle \left(  e_{1}\right\vert
+\left\vert ie_{1}\right\rangle \left(  e_{2}\right\vert
\end{equation}%
\begin{equation}
\mathfrak{C}_{e}^{-1}=\left\vert e_{1}\right)  \left(  e_{1}\right\vert
_{V_{2\mathbb{C}}}+\left\vert e_{2}\right)  \left(  ie_{1}\right\vert
_{V_{2\mathbb{C}}}=\left\vert e_{1}\right)  \left(  e_{1}\right\vert
_{V_{2\mathbb{C}}}+J_{e}\left\vert e_{1}\right)  \left(  ie_{1}\right\vert
_{V_{2\mathbb{C}}}.
\end{equation}


$\left(
\begin{array}
[c]{cc}%
\alpha & -\beta\\
\beta & \alpha
\end{array}
\right)  \in\mathcal{B}\left(  \mathbb{R}^{2}\right)  $ is a representation of
$T_{C}\in\mathcal{B}\left(  \mathbb{C}^{1}\right)  .$ Hence the set of complex
linear transformations $\mathcal{B}\left(  \mathbb{C}^{1}\right)  $ is a real
subalgebra of $\mathcal{B}\left(  \mathbb{R}^{2}\right)  .$

Confirmed by
\begin{equation}
\left(
\begin{array}
[c]{cc}%
t_{R} & -t_{I}\\
t_{I} & t_{R}%
\end{array}
\right)  \left(
\begin{array}
[c]{cc}%
s_{R} & -s_{I}\\
s_{I} & s_{R}%
\end{array}
\right)  =\left(
\begin{array}
[c]{cc}%
t_{R}s_{R}-t_{I}s_{I} & -\left(  t_{R}s_{I}+t_{I}s_{R}\right) \\
\left(  t_{R}s_{I}+t_{I}s_{R}\right)  & t_{R}s_{R}-t_{I}s_{I}%
\end{array}
\right)  .
\end{equation}
Also note
\begin{align}
\left(
\begin{array}
[c]{cc}%
t_{R} & -t_{I}\\
t_{I} & t_{R}%
\end{array}
\right)  \left(
\begin{array}
[c]{cc}%
0 & -1\\
1 & 0
\end{array}
\right)   &  =\left(
\begin{array}
[c]{cc}%
-t_{I} & -t_{R}\\
t_{R} & -t_{I}%
\end{array}
\right) \\
\left(
\begin{array}
[c]{cc}%
0 & -1\\
1 & 0
\end{array}
\right)  \left(
\begin{array}
[c]{cc}%
t_{R} & -t_{I}\\
t_{I} & t_{R}%
\end{array}
\right)   &  =\left(
\begin{array}
[c]{cc}%
-t_{I} & -t_{R}\\
t_{R} & -t_{I}%
\end{array}
\right)
\end{align}
hence
\begin{equation}
\left(
\begin{array}
[c]{cc}%
t_{R} & -t_{I}\\
t_{I} & t_{R}%
\end{array}
\right)  \left(
\begin{array}
[c]{cc}%
0 & -1\\
1 & 0
\end{array}
\right)  =\left(
\begin{array}
[c]{cc}%
0 & -1\\
1 & 0
\end{array}
\right)  \left(
\begin{array}
[c]{cc}%
t_{R} & -t_{I}\\
t_{I} & t_{R}%
\end{array}
\right)
\end{equation}
i.e.
\begin{equation}
\left[  \left(
\begin{array}
[c]{cc}%
t_{R} & -t_{I}\\
t_{I} & t_{R}%
\end{array}
\right)  ,\left(
\begin{array}
[c]{cc}%
0 & -1\\
1 & 0
\end{array}
\right)  \right]  =0.
\end{equation}
If
\begin{equation}
\left(
\begin{array}
[c]{cc}%
a & b\\
c & d
\end{array}
\right)  \left(
\begin{array}
[c]{cc}%
0 & -1\\
1 & 0
\end{array}
\right)  =\left(
\begin{array}
[c]{cc}%
0 & -1\\
1 & 0
\end{array}
\right)  \left(
\begin{array}
[c]{cc}%
a & b\\
c & d
\end{array}
\right)
\end{equation}
then
\begin{align}
\left(
\begin{array}
[c]{cc}%
b & -a\\
d & -c
\end{array}
\right)   &  =\left(
\begin{array}
[c]{cc}%
-c & -d\\
a & b
\end{array}
\right) \\
&  \Rightarrow b=-c,a=d
\end{align}
hence
\begin{equation}
\left(
\begin{array}
[c]{cc}%
a & b\\
c & d
\end{array}
\right)  =\left(
\begin{array}
[c]{cc}%
a & -c\\
c & a
\end{array}
\right)  .
\end{equation}


The subalgebra $\widehat{\mathfrak{C}_{e}}^{-1}\left(  \mathcal{B}\left(
\mathbb{C}^{1}\right)  \right)  \subset\mathcal{B}\left(  \mathbb{R}%
^{2}\right)  $ is thus given by
\begin{equation}
\widehat{\mathfrak{C}_{e}}^{-1}\left(  \mathcal{B}\left(  \mathbb{C}%
^{1}\right)  \right)  =\left\{  \left.  T_{R}\right\vert \left[  T_{R}%
,J_{e}\right]  =0;T_{R}\in\mathcal{B}\left(  \mathbb{R}^{2}\right)  \right\}
\end{equation}
and
\begin{equation}
T_{R}=t_{11}\left\vert e_{1}\right)  \left(  e_{1}\right\vert -t_{12}%
\left\vert e_{1}\right)  \left(  e_{2}\right\vert +t_{12}\left\vert
e_{2}\right)  \left(  e_{1}\right\vert +t_{11}\left\vert e_{2}\right)  \left(
e_{2}\right\vert .
\end{equation}


\paragraph{$\left(  \text{Real}\right)  \dim\left\{  V_{4}\right\}  =4$}

Consider the real space $V_{4},$where
\begin{equation}
\mathbb{R-}\dim\left\{  V_{4}\right\}  =4.
\end{equation}
The basis is
\begin{align}
\left\vert e_{1}\right)   &  =\left\vert \mathbf{e}\right)  \left(
\begin{array}
[c]{c}%
1\\
0\\
0\\
0
\end{array}
\right)  ,\left\vert e_{2}\right)  =\left\vert \mathbf{e}\right)  \left(
\begin{array}
[c]{c}%
0\\
1\\
0\\
0
\end{array}
\right) \\
\left\vert e_{3}\right)   &  =\left\vert \mathbf{e}\right)  \left(
\begin{array}
[c]{c}%
0\\
0\\
1\\
0
\end{array}
\right)  ,\left\vert e_{4}\right)  =\left\vert \mathbf{e}\right)  \left(
\begin{array}
[c]{c}%
0\\
0\\
0\\
1
\end{array}
\right)
\end{align}
and the complex unit $J_{e}$ is given by
\begin{align}
J_{e}\left\vert e_{1}\right)   &  =\left\vert \mathbf{e}\right)  \left(
\begin{array}
[c]{cccc}%
0 & 0 & -1 & 0\\
0 & 0 & 0 & -1\\
1 & 0 & 0 & 0\\
0 & 1 & 0 & 0
\end{array}
\right)  \left(
\begin{array}
[c]{c}%
1\\
0\\
0\\
0
\end{array}
\right)  =\left\vert \mathbf{e}\right)  \left(
\begin{array}
[c]{c}%
0\\
0\\
1\\
0
\end{array}
\right)  =\left\vert e_{3}\right) \\
J_{e}\left\vert e_{2}\right)   &  =\left\vert \mathbf{e}\right)  \left(
\begin{array}
[c]{cccc}%
0 & 0 & -1 & 0\\
0 & 0 & 0 & -1\\
1 & 0 & 0 & 0\\
0 & 1 & 0 & 0
\end{array}
\right)  \left(
\begin{array}
[c]{c}%
0\\
1\\
0\\
0
\end{array}
\right)  =\left\vert \mathbf{e}\right)  \left(
\begin{array}
[c]{c}%
0\\
0\\
0\\
1
\end{array}
\right)  =\left\vert e_{4}\right)
\end{align}
%

\begin{align}
J_{e}\left\vert e_{3}\right)   &  =\left\vert \mathbf{e}\right)  \left(
\begin{array}
[c]{cccc}%
0 & 0 & -1 & 0\\
0 & 0 & 0 & -1\\
1 & 0 & 0 & 0\\
0 & 1 & 0 & 0
\end{array}
\right)  \left(
\begin{array}
[c]{c}%
0\\
0\\
1\\
0
\end{array}
\right)  =\left\vert \mathbf{e}\right)  \left(
\begin{array}
[c]{c}%
-1\\
0\\
0\\
0
\end{array}
\right)  =-\left\vert e_{1}\right)  \\
J_{e}\left\vert e_{4}\right)   &  =\left\vert \mathbf{e}\right)  \left(
\begin{array}
[c]{cccc}%
0 & 0 & -1 & 0\\
0 & 0 & 0 & -1\\
1 & 0 & 0 & 0\\
0 & 1 & 0 & 0
\end{array}
\right)  \left(
\begin{array}
[c]{c}%
0\\
0\\
0\\
1
\end{array}
\right)  =\left\vert \mathbf{e}\right)  \left(
\begin{array}
[c]{c}%
0\\
-1\\
0\\
0
\end{array}
\right)  =-\left\vert e_{2}\right)  ,
\end{align}
i.e.
\begin{equation}
R_{e}\left(  J_{e}\right)  =\left(
\begin{array}
[c]{cc}%
\mathbb{O} & -\mathbb{I}_{2}\\
\mathbb{I}_{2} & \mathbb{O}%
\end{array}
\right)  .
\end{equation}
Hence
\begin{equation}
J_{e}=\left\vert e_{3}\right)  \left(  e_{1}\right\vert +\left\vert
e_{4}\right)  \left(  e_{2}\right\vert -\left\vert e_{1}\right)  \left(
e_{3}\right\vert -\left\vert e_{2}\right)  \left(  e_{4}\right\vert ,
\end{equation}
so corresponds to a real antisymmetric bilinear form, $g,$%
\begin{align}
g\left(  a,b\right)   &  =\left(  \left.  g\right\vert a\wedge b\right)
=\left(  \left.  a\right\vert J_{e}^{t}b\right)  \\
\left\vert g\right)   &  =\left\vert e_{3}\wedge e_{1}\right)  +\left\vert
e_{4}\wedge e_{2}\right)  =\frac{1}{\sqrt{2}}\left\{  \left\vert e_{3}%
e_{1}\right)  +\left\vert e_{4}e_{2}\right)  \right\}  ,
\end{align}
In particular
\begin{equation}
g\left(  e_{j},e_{k}\right)  =\left(  \left.  e_{j}\right\vert J_{e}%
e_{k}\right)  =\left(  \left.  g\right\vert e_{j}\wedge e_{k}\right)
\end{equation}
and%
\begin{equation}
\left\Vert g\right\Vert =1.
\end{equation}
As%
\begin{equation}
\left\vert e_{j}e_{k}\right\rangle \overset{}{=}\sqrt{2}\left\vert e_{j}\wedge
e_{k}\right\rangle
\end{equation}


then%
\begin{equation}
\left(  \left.  g\right\vert e_{1}\wedge e_{3}\right)  =\frac{1}{\sqrt{2}%
}\left(  \left.  g\right\vert e_{1}e_{3}\right)  =\frac{1}{\sqrt{2}}\left(
\left.  \frac{1}{\sqrt{2}}\left\{  \left\vert e_{3}e_{1}\right)  +\left\vert
e_{4}e_{2}\right)  \right\}  \right\vert e_{1}e_{3}\right)  =-\frac{1}{2}%
\end{equation}


$\lceil$%

\begin{align}
\left(  \left.  a\right\vert J_{e}b\right)   &  =\left(
\begin{array}
[c]{cccc}%
a_{1} & a_{2} & a_{3} & a_{4}%
\end{array}
\right)  \left(
\begin{array}
[c]{cccc}%
0 & 0 & -1 & 0\\
0 & 0 & 0 & -1\\
1 & 0 & 0 & 0\\
0 & 1 & 0 & 0
\end{array}
\right)  \left(
\begin{array}
[c]{c}%
b_{1}\\
b_{2}\\
b_{3}\\
b_{4}%
\end{array}
\right) \\
&  =\left(
\begin{array}
[c]{cccc}%
a_{3} & a_{4} & -a_{1} & -a_{2}%
\end{array}
\right)  \left(
\begin{array}
[c]{c}%
b_{1}\\
b_{2}\\
b_{3}\\
b_{4}%
\end{array}
\right)  =a_{3}b_{1}-a_{1}b_{3}+a_{4}b_{2}-a_{2}b_{4}\\
\left(  \left.  g\right\vert a\wedge b\right)   &  =\left(
\begin{array}
[c]{cccccc}%
0 & -1 & 0 & 0 & -1 & 0
\end{array}
\right)  \left(
\begin{array}
[c]{c}%
a_{1}b_{2}-a_{2}b_{1}\\
a_{1}b_{3}-a_{3}b_{1}\\
a_{1}b_{4}-a_{4}b_{1}\\
a_{2}b_{3}-a_{3}b_{2}\\
a_{2}b_{4}-a_{4}b_{2}\\
a_{3}b_{4}-a_{4}b_{3}%
\end{array}
\right)  =a_{3}b_{1}-a_{1}b_{3}+a_{4}b_{2}-a_{2}b_{4}%
\end{align}
$\rfloor$

while
\begin{equation}
I_{4}=\left\vert e_{1}\right)  \left(  e_{1}\right\vert +\left\vert
e_{2}\right)  \left(  e_{2}\right\vert +\left\vert e_{3}\right)  \left(
e_{3}\right\vert +\left\vert e_{4}\right)  \left(  e_{4}\right\vert .
\end{equation}


This real space can be thought of as a complex space of half the dimension
through a complexification map
\begin{equation}
\mathfrak{C}_{e}:V_{4}\rightarrow V_{4e\mathbb{C}},
\end{equation}
given by%
\begin{equation}
\mathfrak{C}_{e}=\left\vert e_{1}\right\rangle \left(  e_{1}\right\vert
+\left\vert e_{2}\right\rangle \left(  e_{2}\right\vert +\left\vert
ie_{3}\right\rangle \left(  e_{3}\right\vert +\left\vert ie_{4}\right\rangle
\left(  e_{4}\right\vert =\left(
\begin{array}
[c]{cccc}%
\left\vert e_{1}\right\rangle  & \left\vert e_{2}\right\rangle  & \left\vert
ie_{1}\right\rangle  & \left\vert ie_{2}\right\rangle
\end{array}
\right)  \left(
\begin{array}
[c]{cccc}%
1 & 0 & 0 & 0\\
0 & 1 & 0 & 0\\
0 & 0 & 1 & 0\\
0 & 0 & 0 & 1
\end{array}
\right)  \left(
\begin{array}
[c]{c}%
\left(  e_{1}\right\vert \\
\left(  e_{2}\right\vert \\
\left(  e_{3}\right\vert \\
\left(  e_{4}\right\vert
\end{array}
\right)  ,
\end{equation}
which has an inverse
\begin{equation}
\mathfrak{C}_{e}^{-1}:V_{4\mathbb{C}}\rightarrow V_{4},
\end{equation}
given by%
\begin{equation}
\mathfrak{C}_{e}^{-1}=\left\vert e_{1}\right)  \left(  e_{1}\right\vert
_{V_{2\mathbb{C}}}+\left\vert e_{2}\right)  \left(  e_{2}\right\vert
_{V_{2\mathbb{C}}}+\left\vert e_{3}\right)  \left(  ie_{1}\right\vert
_{V_{2\mathbb{C}}}+\left\vert e_{4}\right)  \left(  ie_{2}\right\vert
_{V_{2\mathbb{C}}}=\left(
\begin{array}
[c]{cccc}%
\left\vert e_{1}\right)  & \left\vert e_{2}\right)  & \left\vert e_{3}\right)
& \left\vert e_{4}\right)
\end{array}
\right)  \left(
\begin{array}
[c]{cccc}%
1 & 0 & 0 & 0\\
0 & 1 & 0 & 0\\
0 & 0 & 1 & 0\\
0 & 0 & 0 & 1
\end{array}
\right)  \left(
\begin{array}
[c]{c}%
\left(  e_{1}\right\vert _{V_{2\mathbb{C}}}\\
\left(  e_{2}\right\vert _{V_{2\mathbb{C}}}\\
\left(  ie_{1}\right\vert _{V_{2\mathbb{C}}}\\
\left(  ie_{2}\right\vert _{V_{2\mathbb{C}}}%
\end{array}
\right)  .
\end{equation}
The action of $\mathfrak{C}_{e}$ is%
\begin{align}
\mathfrak{C}_{e}\left(  \left\vert a\right)  \right)   &  =\mathfrak{C}%
_{e}\left(  \left\vert e_{1}\right)  a_{11}+\left\vert e_{3}\right)
a_{12}+\left\vert e_{2}\right)  a_{21}+\left\vert e_{4}\right)  a_{22}\right)
=\left\vert e_{1}\right\rangle a_{11}+\left\vert ie_{1}\right\rangle
a_{12}+\left\vert e_{2}\right\rangle a_{21}+\left\vert ie_{2}\right\rangle
a_{22}\\
&  =\left\vert e_{1}\right\rangle a_{11}+\left\vert e_{1}\right\rangle
ia_{12}+\left\vert e_{2}\right\rangle a_{21}+\left\vert e_{2}\right\rangle
ia_{22}=\left\vert e_{1}\right\rangle z_{1}+\left\vert e_{2}\right\rangle
z_{2}=\left\vert e_{1},e_{2}\right\rangle R_{e}\left(  z\right)  ,
\end{align}
which can also be expressed as
\begin{align}
\mathfrak{C}_{e}\left(  \left\vert a\right)  \right)   &  =\mathfrak{C}%
_{e}\left(  \left\vert e_{1}\right)  a_{11}+\mathfrak{C}_{e}\left(  \left\vert
e_{3}\right)  \right)  a_{12}+\mathfrak{C}_{e}\left(  \left\vert e_{2}\right)
\right)  a_{21}+\mathfrak{C}_{e}\left(  \left\vert e_{4}\right)  \right)
a_{22}\right)  =\\
\mathfrak{C}_{e}\left(  \left(  a_{11}I_{4}+a_{12}J_{e}\right)  \left\vert
e_{1}\right)  +\left(  a_{21}I_{4}+a_{22}J_{e}\right)  \left\vert
e_{2}\right)  \right)   &  =\left(  a_{11}+a_{12}i\right)  \left\vert
e_{1}\right\rangle +\left(  a_{21}+a_{22}i\right)  \left\vert e_{2}%
\right\rangle =\left\vert e_{1}\right\rangle z_{1}+\left\vert e_{2}%
\right\rangle z_{2},
\end{align}
i.e.
\begin{equation}
\mathfrak{C}_{e}\left(  \left\vert a\right)  \right)  =\mathfrak{C}_{e}\left(
K_{e}\left(  z_{1}\right)  \left\vert e_{1}\right)  +K_{e}\left(
z_{2}\right)  \left\vert e_{2}\right)  \right)  =\left\vert z_{a}\right\rangle
=z_{1}\left\vert e_{1}\right\rangle +z_{2}\left\vert e_{2}\right\rangle
=\left\vert e_{1},e_{2}\right\rangle R_{e}\left(  z_{a}\right)  ,
\end{equation}
where
\begin{align}
K_{e}\left(  z_{1}\right)   &  =a_{11}R_{e}\left(  I_{4}\right)  +a_{12}%
R_{e}\left(  J_{e}\right) \\
K_{e}\left(  z_{2}\right)   &  =a_{21}R_{e}\left(  I_{4}\right)  +a_{22}%
R_{e}\left(  J_{e}\right)  .
\end{align}


Thus one has
\begin{equation}
\mathfrak{C}_{e}\left(  \left\vert a\right)  \right)  =\left\vert
z\right\rangle =\left\vert z_{1}e_{1}+z_{2}e_{2}\right\rangle =\left\vert
e_{1},e_{2}\right\rangle R_{e}\left(  z\right)  =\mathfrak{C}_{e}\left(
\left\vert I_{4}e_{1},J_{e}e_{1},I_{4}e_{2},J_{e}e_{2}\right)  \left(
\begin{array}
[c]{c}%
a_{11}\\
a_{12}\\
a_{21}\\
a_{22}%
\end{array}
\right)  \right)  =\mathfrak{C}_{e}\left(  \left\vert e_{1},e_{2},e_{3}%
,e_{4}\right)  R_{e}\left(  a\right)  \right)  ,
\end{equation}
where
\begin{equation}
K_{e}\left(  z_{j}\right)  =\left(  a_{j}I_{4}+a_{j+2}J_{e}\right)
\end{equation}
is a representation of $\mathbb{C}$ in $\mathcal{B}\left(  V_{4}\right)  .$

\subsubsection{\bigskip Linear Transformation}

The relationship between complex linear transformation
\begin{equation}
T_{C}:V_{2e\mathbb{C}}\rightarrow V_{2e\mathbb{C}}%
\end{equation}
and real linear transformation
\begin{equation}
T_{R}:V_{2}\rightarrow V_{2}%
\end{equation}
is given by%
\begin{equation}
T_{C}=\mathfrak{C}_{e}T_{R}\mathfrak{C}_{e}^{-1}=\widehat{\mathfrak{C}}%
_{e}\left(  T_{R}\right)
\end{equation}
The action of $T_{C}$ is given by
\begin{align}
\left\vert w\right\rangle  &  =T_{C}\left\vert z\right\rangle =T_{C}\left\vert
e_{1},e_{2}\right\rangle R_{e}\left(  z\right)  =\left\vert e_{1}%
,e_{2}\right\rangle R_{e}\left(  T_{\mathbb{C}}\right)  R_{e}\left(  z\right)
\\
&  =\left\vert e_{1},e_{2}\right\rangle \left(
\begin{array}
[c]{cc}%
T_{11} & T_{12}\\
T_{21} & T_{22}%
\end{array}
\right)  \left(
\begin{array}
[c]{c}%
z_{1}\\
z_{2}%
\end{array}
\right)  =\left\vert e_{1},e_{2}\right\rangle \left(
\begin{array}
[c]{cc}%
x_{11}+iy_{11} & x_{12}+iy_{12}\\
x_{21}+iy_{21} & x_{22}+iy_{22}%
\end{array}
\right)  \left(
\begin{array}
[c]{c}%
a_{11}+ia_{12}\\
a_{21}+ia_{22}%
\end{array}
\right) \\
&  =\left\vert e_{1},ie_{1},e_{2},ie_{2}\right\rangle \left(
\begin{array}
[c]{cc}%
\mathbb{T}_{11} & \mathbb{T}_{12}\\
\mathbb{T}_{21} & \mathbb{T}_{22}%
\end{array}
\right)  \left(
\begin{array}
[c]{c}%
a_{11}\\
a_{12}\\
a_{21}\\
a_{22}%
\end{array}
\right)  =\left\vert e_{1},ie_{1},e_{2},ie_{2}\right\rangle \left(
\begin{array}
[c]{cc}%
\left(
\begin{array}
[c]{cc}%
x_{11} & -y_{11}\\
y_{11} & x_{11}%
\end{array}
\right)  & \left(
\begin{array}
[c]{cc}%
x_{12} & -y_{12}\\
y_{12} & x_{12}%
\end{array}
\right) \\
\left(
\begin{array}
[c]{cc}%
x_{21} & -y_{21}\\
y_{21} & x_{21}%
\end{array}
\right)  & \left(
\begin{array}
[c]{cc}%
x_{22} & -y_{22}\\
y_{22} & x_{22}%
\end{array}
\right)
\end{array}
\right)  \left(
\begin{array}
[c]{c}%
a_{11}\\
a_{12}\\
a_{21}\\
a_{22}%
\end{array}
\right) \\
&  =\left\vert e_{1},e_{2},ie_{1},ie_{2}\right\rangle \left(
\begin{array}
[c]{cc}%
\mathbb{T}_{R} & -\mathbb{T}_{I}\\
\mathbb{T}_{I} & \mathbb{T}_{R}%
\end{array}
\right)  \left(
\begin{array}
[c]{c}%
a_{11}\\
a_{21}\\
a_{12}\\
a_{22}%
\end{array}
\right)  =\left\vert e_{1},e_{2},ie_{1},ie_{2}\right\rangle \left(
\begin{array}
[c]{cc}%
\left(
\begin{array}
[c]{cc}%
x_{11} & x_{12}\\
x_{21} & x_{22}%
\end{array}
\right)  & -\left(
\begin{array}
[c]{cc}%
y_{11} & y_{12}\\
y_{21} & y_{22}%
\end{array}
\right) \\
\left(
\begin{array}
[c]{cc}%
y_{11} & y_{12}\\
y_{21} & y_{22}%
\end{array}
\right)  & \left(
\begin{array}
[c]{cc}%
x_{11} & x_{12}\\
x_{21} & x_{22}%
\end{array}
\right)
\end{array}
\right)  \left(
\begin{array}
[c]{c}%
a_{11}\\
a_{21}\\
a_{12}\\
a_{22}%
\end{array}
\right)
\end{align}


The subalgebra $\widehat{\mathfrak{C}_{e}}^{-1}\left(  \mathcal{B}\left(
\mathbb{C}^{2}\right)  \right)  \subset\mathcal{B}\left(  \mathbb{R}%
^{4}\right)  $ is thus given by
\begin{equation}
\widehat{\mathfrak{C}_{e}}^{-1}\left(  \mathcal{B}\left(  \mathbb{C}%
^{2}\right)  \right)  =\left\{  \left.  T_{R}\right\vert \left[  T_{R}%
,J_{e}\right]  =0;T_{R}\in\mathcal{B}\left(  \mathbb{R}^{4}\right)  \right\}
,
\end{equation}
and
\begin{align}
T_{R}  &  =t_{11}\left\vert e_{1}\right)  \left(  e_{1}\right\vert
+t_{12}\left\vert e_{1}\right)  \left(  e_{2}\right\vert -t_{13}\left\vert
e_{1}\right)  \left(  e_{3}\right\vert -t_{14}\left\vert e_{1}\right)  \left(
e_{4}\right\vert \\
&  t_{21}\left\vert e_{2}\right)  \left(  e_{1}\right\vert +t_{22}\left\vert
e_{2}\right)  \left(  e_{2}\right\vert -t_{23}\left\vert e_{2}\right)  \left(
e_{3}\right\vert -t_{24}\left\vert e_{2}\right)  \left(  e_{4}\right\vert \\
&  t_{13}\left\vert e_{3}\right)  \left(  e_{1}\right\vert +t_{14}\left\vert
e_{3}\right)  \left(  e_{2}\right\vert +t_{11}\left\vert e_{3}\right)  \left(
e_{3}\right\vert +t_{12}\left\vert e_{3}\right)  \left(  e_{4}\right\vert \\
&  t_{23}\left\vert e_{4}\right)  \left(  e_{1}\right\vert +t_{24}\left\vert
e_{4}\right)  \left(  e_{2}\right\vert +t_{21}\left\vert e_{4}\right)  \left(
e_{3}\right\vert +t_{22}\left\vert e_{4}\right)  \left(  e_{4}\right\vert ,
\end{align}
which shows that $\widehat{\mathfrak{C}_{e}}^{-1}\left(  \mathcal{B}\left(
\mathbb{C}^{2}\right)  \right)  $ is an invariant subspace of the one
parameter group $\widehat{\mathfrak{u}_{e}}$ defined by
\begin{equation}
\widehat{\mathfrak{u}_{e}\left(  x\right)  }\left(  T_{R}\right)  =e^{xJ_{e}%
}T_{R}e^{-xJ_{e}}=T_{R}\,\forall T_{R}\in\widehat{\mathfrak{C}_{e}}%
^{-1}\left(  \mathcal{B}\left(  \mathbb{C}^{2}\right)  \right)  ;x\in
\mathbb{R}.
\end{equation}


The operator $J_{e}$ belongs to the orthogonal group $O\left(  4,\mathbb{R}%
\right)  $ as
\begin{gather}
\left(  J_{e}^{t}=-J_{e}\right)  \,\&\left(  J_{e}^{2}=-I_{4}\right) \\
\Downarrow\\
J_{e}\in O\left(  \mathbb{R}^{4}\right)
\end{gather}
$\lceil$%

\begin{equation}
\left(  -J_{e}\right)  J_{e}=I_{4}\Rightarrow J_{e}^{-1}=\left(
-J_{e}\right)  =J_{e}^{t}%
\end{equation}
$\rfloor$

and it also belongs to the Lie algebra $o\left(  4,\mathbb{R}\right)  $ as
\begin{equation}
J_{e}^{t}=-J_{e}.
\end{equation}


\begin{theorem}
The set $\widehat{\mathfrak{C}_{e}}^{-1}\left(  \mathcal{B}\left(
\mathbb{C}^{2}\right)  \right)  =\left\{  \left.  T_{R}\right\vert \left[
T_{R},J_{e}\right]  =0;T_{R}\in\mathcal{B}\left(  \mathbb{R}^{4}\right)
\right\}  $ is invariant under the transformations$\left\{
\widehat{\mathcal{C}}\right\}  $ where
\begin{equation}
\widehat{\mathcal{C}}\left(  J_{e}\right)  =\mathcal{C}J_{e}\mathcal{C}%
^{t}=J_{e},
\end{equation}
which form the adjoint representation of the symplectic group $Sp\left(
4,\mathbb{R}\right)  $ that is defined by
\begin{equation}
Sp\left(  4,\mathbb{R}\right)  =\left\{  \left.  \mathcal{C}\right\vert
\mathcal{C}J_{e}\mathcal{C}^{t}=J_{e};\mathcal{C}\in Gl\left(  4,\mathbb{R}%
\right)  \subset\mathcal{B}\left(  \mathbb{R}^{4}\right)  \right\}  .
\end{equation}

\end{theorem}

$\lceil$Note that if $\mathcal{C\in}Sp\left(  4,\mathbb{R}\right)  $ then
\begin{align}
-J_{e}\mathcal{C}J_{e}\mathcal{C}^{t}  &  =I_{4}\\
\left(  J_{e}\mathcal{C}^{t}\right)  ^{-1}  &  =-J_{e}\mathcal{C}\\
J_{e}\mathcal{C}^{t}J_{e}\mathcal{C}  &  =\mathcal{-}I_{4}\\
\mathcal{C}^{t}J_{e}\mathcal{C}  &  =J_{e}%
\end{align}
i.e.
\begin{equation}
\mathcal{C}J_{e}\mathcal{C}^{t}=J_{e}\Longleftrightarrow\mathcal{C}^{t}%
J_{e}\mathcal{C}=J_{e}.
\end{equation}
$\rfloor$ As
\begin{equation}
\left[  T_{R},J_{e}\right]  =0
\end{equation}
one obtains that
\begin{align}
\mathcal{C}\left[  T_{R},J_{e}\right]  \mathcal{C}^{-1}  &  =\mathcal{C}%
T_{R}^{-1}\mathcal{CC}J_{e}^{-1}\mathcal{C}-\mathcal{C}J_{e}^{-1}%
\mathcal{CC}T_{R}^{-1}\mathcal{C}\\
&  =\left[  \mathcal{C}T_{R}^{-1}\mathcal{C},J_{e}^{\prime}\right]  =0.
\end{align}
If
\begin{equation}
J_{e}^{\prime}=\mathcal{C}J_{e}^{-1}\mathcal{C}=J_{e}%
\end{equation}
i.e. $\mathcal{C}J_{e}^{-1}\mathcal{C}\in\widehat{\mathfrak{C}_{e}}%
^{-1}\left(  \mathcal{B}\left(  \mathbb{C}^{2}\right)  \right)  $ then
\begin{equation}
\mathcal{C}^{-1}=\mathcal{C}^{t}%
\end{equation}
i.e. $\mathcal{C}\in Sp\left(  4,\mathbb{R}\right)  $ and
\begin{equation}
\widehat{\mathcal{C}}:\widehat{\mathfrak{C}_{e}}^{-1}\left(  \mathcal{B}%
\left(  \mathbb{C}^{2}\right)  \right)  \rightarrow\widehat{\mathfrak{C}_{e}%
}^{-1}\left(  \mathcal{B}\left(  \mathbb{C}^{2}\right)  \right)  .
\end{equation}


If $\mathcal{C}\notin Sp\left(  4,\mathbb{R}\right)  $ then
\begin{equation}
\mathcal{C}J_{e}\mathcal{C}^{t}=-\left(  \mathcal{C}J_{e}\mathcal{C}%
^{t}\right)  ^{t}%
\end{equation}
and
\begin{equation}
\mathcal{C}J_{e}\mathcal{C}^{t}\mathcal{C}J_{e}\mathcal{C}^{t}=-I_{4}%
\end{equation}
if and only if
\begin{equation}
\mathcal{C}^{t}\mathcal{C}=\mathcal{CC}^{t}=I_{4}%
\end{equation}
i.e. $O\in O\left(  4,\mathbb{R}\right)  .$

Hence complexifications are related to each by
\begin{align}
J_{e}^{\prime}  &  =OJ_{e}O^{t}=\widehat{O}\left(  J_{e}\right) \\
O  &  \in O\left(  4,\mathbb{R}\right)  .
\end{align}
The set of unique complexifications $\mathfrak{C0m}$ thus corresponds to
\begin{equation}
\Theta\in O\left(  4,\mathbb{R}\right)  ;\Theta\notin Sp\left(  4,\mathbb{R}%
\right)
\end{equation}
i.e.
\begin{equation}
\mathfrak{C0m}=\left.  O\left(  4,\mathbb{R}\right)  \right/  \left(  O\left(
4,\mathbb{R}\right)  \cap Sp\left(  4,\mathbb{R}\right)  \right)  .
\end{equation}


\paragraph{Inner products}

The inner product in $\mathbb{R}^{4}$ induces an inner product in
$\mathbb{C}^{2}$ by
\begin{align}
\left\langle \left.  \omega\right\vert z\right\rangle  &  =\operatorname{Re}%
\left\langle \left.  \omega\right\vert z\right\rangle +i\operatorname{Im}%
\left\langle \left.  \omega\right\vert z\right\rangle =re^{i\theta}\\
&  =\left\langle \left.  \mathfrak{C}_{e}\left(  a\right)  \right\vert
\mathfrak{C}_{e}\left(  b\right)  \right\rangle =\left(  \left.
\mathfrak{C}_{e}^{-1}\left(  \omega\right)  \right\vert \mathfrak{C}_{e}%
^{-1}\left(  z\right)  \right) \\
&  =\left(  \left.  a\right\vert b\right)  +i\left(  \left.  a\right\vert
J_{e}b\right)
\end{align}
and then
\begin{align}
\operatorname{Re}\left\langle \left.  \omega\right\vert z\right\rangle  &
=\left(  \left.  a\right\vert b\right) \\
\operatorname{Im}\left\langle \left.  \omega\right\vert z\right\rangle  &
=\left(  \left.  a\right\vert J_{e}b\right)  .
\end{align}
In terms of representations
\begin{align}
&  R_{e}\left(  \omega\right)  ^{\dagger}\left\langle \left.  e_{1}%
,e_{2}\right\vert e_{1},e_{2}\right\rangle R_{e}\left(  z\right)  =\left(
\begin{array}
[c]{cc}%
\overline{\omega_{1}} & \overline{\omega_{2}}%
\end{array}
\right)  \left(
\begin{array}
[c]{cc}%
\left\langle \left.  e_{1}\right\vert e_{1}\right\rangle  & \left\langle
\left.  e_{1}\right\vert e_{2}\right\rangle \\
\left\langle \left.  e_{2}\right\vert e_{1}\right\rangle  & \left\langle
\left.  e_{2}\right\vert e_{2}\right\rangle
\end{array}
\right)  \left(
\begin{array}
[c]{c}%
z_{1}\\
z_{2}%
\end{array}
\right) \\
&  =\left(
\begin{array}
[c]{cc}%
a_{11}-ia_{12} & a_{21}-ia_{22}%
\end{array}
\right)  \left(
\begin{array}
[c]{cc}%
1 & 0\\
0 & 1
\end{array}
\right)  \left(
\begin{array}
[c]{c}%
b_{11}+ib_{12}\\
b_{21}+ib_{22}%
\end{array}
\right) \\
&  =\left(  \left(
\begin{array}
[c]{cc}%
a_{11} & a_{21}%
\end{array}
\right)  -i\left(
\begin{array}
[c]{cc}%
a_{12} & a_{22}%
\end{array}
\right)  \right)  \left(
\begin{array}
[c]{cc}%
1 & 0\\
0 & 1
\end{array}
\right)  \left(  \left(
\begin{array}
[c]{c}%
b_{11}\\
b_{21}%
\end{array}
\right)  +i\left(
\begin{array}
[c]{c}%
b_{12}\\
b_{22}%
\end{array}
\right)  \right) \\
&  =\left(
\begin{array}
[c]{cc}%
a_{11} & a_{21}%
\end{array}
\right)  \left(
\begin{array}
[c]{c}%
b_{11}\\
b_{21}%
\end{array}
\right)  +\left(
\begin{array}
[c]{cc}%
a_{12} & a_{22}%
\end{array}
\right)  \left(
\begin{array}
[c]{c}%
b_{12}\\
b_{22}%
\end{array}
\right)  +i\left(  \left(
\begin{array}
[c]{cc}%
a_{11} & a_{21}%
\end{array}
\right)  \left(
\begin{array}
[c]{c}%
b_{12}\\
b_{22}%
\end{array}
\right)  -\left(
\begin{array}
[c]{cc}%
a_{12} & a_{22}%
\end{array}
\right)  \left(
\begin{array}
[c]{c}%
b_{11}\\
b_{21}%
\end{array}
\right)  \right) \\
&  =\left(
\begin{array}
[c]{cccc}%
a_{11} & a_{21} & a_{12} & a_{22}%
\end{array}
\right)  \left(
\begin{array}
[c]{cc}%
\left(
\begin{array}
[c]{cc}%
1 & 0\\
0 & 1
\end{array}
\right)  & \left(
\begin{array}
[c]{cc}%
0 & 0\\
0 & 0
\end{array}
\right) \\
\left(
\begin{array}
[c]{cc}%
0 & 0\\
0 & 0
\end{array}
\right)  & \left(
\begin{array}
[c]{cc}%
1 & 0\\
0 & 1
\end{array}
\right)
\end{array}
\right)  \left(
\begin{array}
[c]{c}%
b_{11}\\
b_{21}\\
b_{12}\\
b_{22}%
\end{array}
\right) \\
&  +i\left(
\begin{array}
[c]{cccc}%
a_{11} & a_{21} & a_{12} & a_{22}%
\end{array}
\right)  \left(
\begin{array}
[c]{cc}%
\left(
\begin{array}
[c]{cc}%
0 & 0\\
0 & 0
\end{array}
\right)  & -\left(
\begin{array}
[c]{cc}%
1 & 0\\
0 & 1
\end{array}
\right) \\
\left(
\begin{array}
[c]{cc}%
1 & 0\\
0 & 1
\end{array}
\right)  & \left(
\begin{array}
[c]{cc}%
0 & 0\\
0 & 0
\end{array}
\right)
\end{array}
\right)  \left(
\begin{array}
[c]{c}%
b_{11}\\
b_{21}\\
b_{12}\\
b_{22}%
\end{array}
\right) \\
&  =R_{e}\left(  a\right)  ^{t}S_{e}R_{e}\left(  b\right)  +iR_{e}\left(
a\right)  ^{t}J_{e}R_{e}\left(  b\right)  .
\end{align}


If $T_{R}\in O\left(  \mathbb{R}^{4}\right)  \cap Sp\left(  \mathbb{R}%
^{4}\right)  $ i.e.
\begin{align}
T_{R}^{t}T_{R}  &  =I_{\mathbb{R}^{4}}\\
T_{R}^{t}J_{e}T_{R}  &  =J_{e}%
\end{align}
then
\begin{equation}
\left[  T_{R},J_{e}\right]  =0
\end{equation}
and $T_{R}\in U\left(  \mathbb{C}^{2}\right)  $ i.e.
\begin{equation}
\widehat{\mathfrak{C}_{e}}\left(  T_{R}\right)  ^{\dagger}%
\widehat{\mathfrak{C}_{e}}\left(  T_{R}\right)  =I_{\mathbb{C}^{2}}%
\end{equation}
as
\begin{align}
\left\langle \left.  T_{\mathbb{C}}\omega\right\vert T_{\mathbb{C}%
}z\right\rangle  &  =\operatorname{Re}\left\langle \left.  T_{\mathbb{C}%
}\omega\right\vert T_{\mathbb{C}}z\right\rangle +i\operatorname{Im}%
\left\langle \left.  T_{\mathbb{C}}\omega\right\vert T_{\mathbb{C}%
}z\right\rangle =re^{i\theta}\\
&  =\left(  \left.  T_{\mathbb{R}}a\right\vert T_{\mathbb{R}}b\right)
+i\left(  \left.  T_{\mathbb{R}}a\right\vert J_{e}T_{\mathbb{R}}b\right) \\
&  =\left(  \left.  a\right\vert b\right)  +i\left(  \left.  a\right\vert
J_{e}b\right)  .
\end{align}
This shows that
\begin{equation}
\mathfrak{C0m}=\left.  O\left(  4,\mathbb{R}\right)  \right/  U\left(
2,\mathbb{C}\right)  .
\end{equation}


The different complex algebras represented in $\mathcal{B}\left(
\mathbb{R}^{4}\right)  $ are thus connected by
\begin{equation}
\widehat{\Theta}:\widehat{\mathfrak{C}_{e_{1}}}^{-1}\left(  \mathcal{B}\left(
\mathbb{C}^{2}\right)  \right)  \rightarrow\widehat{\mathfrak{C}_{e_{2}}}%
^{-1}\left(  \mathcal{B}\left(  \mathbb{C}^{2}\right)  \right)
;\widehat{\Theta}\in\mathfrak{C0m.}%
\end{equation}


\subsection{Questions}

\begin{itemize}
\item $J_{e_{1}}\neq J_{e_{2}}$
\begin{equation}
\left[  T_{R},J_{e_{1}}\right]  =0\Rightarrow\left[  T_{R},J_{e_{2}}\right]
\neq0
\end{equation}
%

\begin{align}
T_{R}J_{e_{1}}-J_{e_{1}}T_{R}  &  =0\Leftrightarrow T_{R}=J_{e_{1}}%
T_{R}J_{e_{1}}^{t}\\
T_{R}J_{e_{2}}-J_{e_{2}}T_{R}  &  =0\Leftrightarrow T_{R}=J_{e_{2}}%
T_{R}J_{e_{2}}^{t}\\
&  \Rightarrow J_{e_{2}}T_{R}J_{e_{2}}^{t}=J_{e_{1}}T_{R}J_{e_{1}}^{t}\\
&  \Rightarrow T_{R}=J_{e_{2}}^{t}J_{e_{1}}T_{R}J_{e_{1}}^{t}J_{e_{2}}\\
&  \Rightarrow J_{e_{1}}^{t}J_{e_{2}}T_{R}=T_{R}J_{e_{1}}^{t}J_{e_{2}}%
\end{align}%
\begin{align}
J_{e_{2}}^{t}J_{e_{1}}  &  =J_{e_{2}}J_{e_{1}}^{t}\\
\left(  J_{e_{2}}^{t}J_{e_{1}}\right)  ^{t}  &  =J_{e_{1}}^{t}J_{e_{2}%
}=J_{e_{1}}J_{e_{2}}^{t}%
\end{align}


\item
\begin{equation}
\exists J_{e}\forall T_{R}\in\mathcal{B}\left(  \mathbb{R}^{4}\right)
\end{equation}


\item
\begin{gather}
\mathcal{B}\left(  \mathbb{R}^{4}\right) \\
\widehat{\mathfrak{C}_{e_{1}}}\swarrow\cdots\searrow\widehat{\mathfrak{C}%
_{e_{2}}}\\
\mathcal{B}\left(  \mathbb{C}^{2}\right)  _{e_{1}}\overset{u_{12}%
}{\longrightarrow}\mathcal{B}\left(  \mathbb{C}^{2}\right)  _{e_{2}}%
\end{gather}


Complex units are related to each other by
\begin{equation}
J_{e_{2}}=\widehat{O_{12}}\left(  J_{e_{1}}\right)  =O_{12}J_{e_{1}}O_{12}%
^{t};O_{12}\in O\left(  4,\mathbb{R}\right)  ,
\end{equation}
which clearly shows that if $O_{12}\in O\left(  4,\mathbb{R}\right)  \cap
Sp\left(  4,\mathbb{R}\right)  $ then $J_{e_{2}}=J_{e_{1}}.$
\end{itemize}

%

\begin{align}
T_{R}  &  =t_{11}\left\vert e_{1}\right)  \left(  e_{1}\right\vert
+t_{12}\left\vert e_{1}\right)  \left(  e_{2}\right\vert -t_{13}\left\vert
e_{1}\right)  \left(  e_{3}\right\vert -t_{14}\left\vert e_{1}\right)  \left(
e_{4}\right\vert \\
&  t_{21}\left\vert e_{2}\right)  \left(  e_{1}\right\vert +t_{22}\left\vert
e_{2}\right)  \left(  e_{2}\right\vert -t_{23}\left\vert e_{2}\right)  \left(
e_{3}\right\vert -t_{24}\left\vert e_{2}\right)  \left(  e_{4}\right\vert \\
&  t_{13}\left\vert e_{3}\right)  \left(  e_{1}\right\vert +t_{14}\left\vert
e_{3}\right)  \left(  e_{2}\right\vert +t_{11}\left\vert e_{3}\right)  \left(
e_{3}\right\vert +t_{12}\left\vert e_{3}\right)  \left(  e_{4}\right\vert \\
&  t_{23}\left\vert e_{4}\right)  \left(  e_{1}\right\vert +t_{24}\left\vert
e_{4}\right)  \left(  e_{2}\right\vert +t_{21}\left\vert e_{4}\right)  \left(
e_{3}\right\vert +t_{22}\left\vert e_{4}\right)  \left(  e_{4}\right\vert ,
\end{align}%
\begin{equation}
T_{R}=\left(
\begin{array}
[c]{cc}%
\left\vert \mathbf{e}_{R}\right)  & \left\vert \mathbf{e}_{I}\right)
\end{array}
\right)  \left(
\begin{array}
[c]{cc}%
t_{R} & -t_{I}\\
t_{I} & t_{R}%
\end{array}
\right)  \left(
\begin{array}
[c]{c}%
\left(  \mathbf{e}_{R}\right\vert \\
\left(  \mathbf{e}_{I}\right\vert
\end{array}
\right)
\end{equation}%
\begin{align}
T_{R}  &  =s_{11}\left\vert f_{1}\right)  \left(  f_{1}\right\vert
+s_{12}\left\vert f_{1}\right)  \left(  f_{2}\right\vert -s_{13}\left\vert
f_{1}\right)  \left(  f_{3}\right\vert -s_{14}\left\vert f_{1}\right)  \left(
f_{4}\right\vert \\
&  s_{21}\left\vert f_{2}\right)  \left(  f_{1}\right\vert +s_{22}\left\vert
f_{2}\right)  \left(  f_{2}\right\vert -s_{23}\left\vert f_{2}\right)  \left(
f_{3}\right\vert -s_{24}\left\vert f_{2}\right)  \left(  f_{4}\right\vert \\
&  s_{13}\left\vert f_{3}\right)  \left(  f_{1}\right\vert +s_{14}\left\vert
f_{3}\right)  \left(  f_{2}\right\vert +s_{11}\left\vert f_{3}\right)  \left(
f_{3}\right\vert +s_{12}\left\vert f_{3}\right)  \left(  f_{4}\right\vert \\
&  s_{23}\left\vert f_{4}\right)  \left(  f_{1}\right\vert +s_{24}\left\vert
f_{4}\right)  \left(  f_{2}\right\vert +s_{21}\left\vert f_{4}\right)  \left(
f_{3}\right\vert +s_{22}\left\vert f_{4}\right)  \left(  f_{4}\right\vert ,
\end{align}%
\begin{equation}
T_{R}=\left(
\begin{array}
[c]{cc}%
\left\vert \mathbf{f}_{R}\right)  & \left\vert \mathbf{f}_{I}\right)
\end{array}
\right)  \left(
\begin{array}
[c]{cc}%
s_{R} & -s_{I}\\
s_{I} & s_{R}%
\end{array}
\right)  \left(
\begin{array}
[c]{c}%
\left(  \mathbf{f}_{R}\right\vert \\
\left(  \mathbf{f}_{I}\right\vert
\end{array}
\right)  =\left(
\begin{array}
[c]{cc}%
\left\vert \mathbf{e}_{R}\right)  & \left\vert \mathbf{e}_{I}\right)
\end{array}
\right)  R_{e}\left(  O_{12}^{t}\right)  \left(
\begin{array}
[c]{cc}%
s_{R} & -s_{I}\\
s_{I} & s_{R}%
\end{array}
\right)  R_{e}\left(  O_{12}\right)  \left(
\begin{array}
[c]{c}%
\left(  \mathbf{e}_{R}\right\vert \\
\left(  \mathbf{e}_{I}\right\vert
\end{array}
\right)
\end{equation}%
\begin{equation}
\left(
\begin{array}
[c]{cc}%
t_{R} & -t_{I}\\
t_{I} & t_{R}%
\end{array}
\right)  =R_{e}\left(  O_{12}^{t}\right)  \left(
\begin{array}
[c]{cc}%
s_{R} & -s_{I}\\
s_{I} & s_{R}%
\end{array}
\right)  R_{e}\left(  O_{12}\right)
\end{equation}
Hence
\begin{align}
&  T_{R}J_{e_{2}}-J_{e_{2}}T_{R}=0\\
&  \Rightarrow T_{R}O_{12}J_{e_{1}}O_{12}^{t}-O_{12}J_{e_{1}}O_{12}^{t}%
T_{R}=0\\
&  \Rightarrow O_{12}^{t}T_{R}O_{12}J_{e_{1}}-J_{e_{1}}O_{12}^{t}T_{R}O_{12}=0
\end{align}
i.e.%
\begin{equation}
O_{12}^{t}T_{R}O_{12}\in\widehat{\mathfrak{C}_{e_{1}}}^{-1}\left(
\mathcal{B}\left(  \mathbb{C}^{2}\right)  \right)  .
\end{equation}
So if
\begin{equation}
O_{12}^{t}T_{R}O_{12}\text{ and }T_{R}\in\widehat{\mathfrak{C}_{e_{1}}}%
^{-1}\left(  \mathcal{B}\left(  \mathbb{C}^{2}\right)  \right)
\end{equation}
then
\begin{align}
T_{R}J_{e_{1}}-J_{e_{1}}T_{R}  &  =0\Rightarrow T_{R}=J_{e_{1}}^{t}%
T_{R}J_{e_{1}}\\
O_{12}^{t}T_{R}O_{12}J_{e_{1}}-J_{e_{1}}O_{12}^{t}T_{R}O_{12}  &
=0\Rightarrow T_{R}=O_{12}J_{e_{1}}O_{12}^{t}T_{R}O_{12}J_{e_{1}}^{t}%
O_{12}^{t}.
\end{align}


$\lceil$%
\begin{equation}
T_{R}J_{e_{2}}-J_{e_{2}}T_{R}=0\Leftrightarrow J_{e_{2}}^{t}T_{R}^{t}%
-T_{R}^{t}J_{e_{2}}^{t}=T_{R}^{t}J_{e_{2}}-J_{e_{2}}T_{R}^{t}=\left[
T_{R}^{t},J_{e_{2}}\right]
\end{equation}
$\rfloor$

\subsection{Groups}

Let%
\begin{equation}
Q\left(
\begin{array}
[c]{cc}%
\left\vert \mathbf{e}_{R}\right)  & \left\vert \mathbf{e}_{I}\right)
\end{array}
\right)  =\left(
\begin{array}
[c]{cc}%
\left\vert \mathbf{e}_{R}\right)  & \left\vert \mathbf{e}_{I}\right)
\end{array}
\right)  \left(
\begin{array}
[c]{cc}%
Q_{RR} & Q_{RI}\\
Q_{IR} & Q_{II}%
\end{array}
\right)
\end{equation}
then%
\begin{align}
Q  &  \in O\left(  4,\mathbb{R}\right)  \Leftrightarrow\left(
\begin{array}
[c]{cc}%
Q_{RR} & Q_{RI}\\
Q_{IR} & Q_{II}%
\end{array}
\right)  \left(
\begin{array}
[c]{cc}%
Q_{RR} & Q_{RI}\\
Q_{IR} & Q_{II}%
\end{array}
\right)  ^{t}=\left(
\begin{array}
[c]{cc}%
Q_{RR} & Q_{RI}\\
Q_{IR} & Q_{II}%
\end{array}
\right)  \left(
\begin{array}
[c]{cc}%
Q_{RR}^{t} & Q_{IR}^{t}\\
Q_{RI}^{t} & Q_{II}^{t}%
\end{array}
\right) \\
&  =\left(
\begin{array}
[c]{cc}%
\mathbb{I}_{2} & \mathbb{O}\\
\mathbb{O} & \mathbb{I}_{2}%
\end{array}
\right)  =\left(
\begin{array}
[c]{cc}%
Q_{RR}Q_{RR}^{t}+Q_{RI}Q_{RI}^{t} & Q_{RR}Q_{IR}^{t}+Q_{RI}Q_{II}^{t}\\
Q_{IR}Q_{RR}^{t}+Q_{II}Q_{RI}^{t} & Q_{IR}Q_{IR}^{t}+Q_{II}Q_{II}^{t}%
\end{array}
\right) \\
&  Q_{RR}Q_{RR}^{t}+Q_{RI}Q_{RI}^{t}=\mathbb{I}_{2}\\
&  Q_{IR}Q_{IR}^{t}+Q_{II}Q_{II}^{t}=\mathbb{I}_{2}\\
&  Q_{IR}Q_{RR}^{t}+Q_{II}Q_{RI}^{t}=\mathbb{O}\\
&  Q_{RR}Q_{IR}^{t}+Q_{RI}Q_{II}^{t}=\mathbb{O}%
\end{align}
and%
\begin{align}
&  Q\in Sp\left(  4,\mathbb{R}\right)  \Leftrightarrow\left(
\begin{array}
[c]{cc}%
Q_{RR} & Q_{RI}\\
Q_{IR} & Q_{II}%
\end{array}
\right)  \left(
\begin{array}
[c]{cc}%
\mathbb{O} & \mathbb{-I}_{2}\\
\mathbb{I}_{2} & \mathbb{O}%
\end{array}
\right)  \left(
\begin{array}
[c]{cc}%
Q_{RR} & Q_{RI}\\
Q_{IR} & Q_{II}%
\end{array}
\right)  ^{t}\\
&  =\left(
\begin{array}
[c]{cc}%
Q_{RR} & Q_{RI}\\
Q_{IR} & Q_{II}%
\end{array}
\right)  \left(
\begin{array}
[c]{cc}%
\mathbb{O} & \mathbb{-I}_{2}\\
\mathbb{I}_{2} & \mathbb{O}%
\end{array}
\right)  \left(
\begin{array}
[c]{cc}%
Q_{RR}^{t} & Q_{IR}^{t}\\
Q_{RI}^{t} & Q_{II}^{t}%
\end{array}
\right) \\
&  =\left(
\begin{array}
[c]{cc}%
\mathbb{O} & \mathbb{-I}_{2}\\
\mathbb{I}_{2} & \mathbb{O}%
\end{array}
\right)  =\left(
\begin{array}
[c]{cc}%
Q_{RI} & -Q_{RR}\\
Q_{II} & -Q_{IR}%
\end{array}
\right)  \left(
\begin{array}
[c]{cc}%
Q_{RR}^{t} & Q_{IR}^{t}\\
Q_{RI}^{t} & Q_{II}^{t}%
\end{array}
\right)  =\left(
\begin{array}
[c]{cc}%
Q_{RI}Q_{RR}^{t}-Q_{RR}Q_{RI}^{t} & Q_{RI}Q_{IR}^{t}-Q_{RR}Q_{II}^{t}\\
Q_{II}Q_{RR}^{t}-Q_{IR}Q_{RI}^{t} & Q_{II}Q_{IR}^{t}-Q_{IR}Q_{II}^{t}%
\end{array}
\right) \\
&  Q_{RI}Q_{RR}^{t}-Q_{RR}Q_{RI}^{t}=\mathbb{O}\\
&  Q_{II}Q_{IR}^{t}-Q_{IR}Q_{II}^{t}=\mathbb{O}\\
&  Q_{II}Q_{RR}^{t}-Q_{IR}Q_{RI}^{t}=\mathbb{I}_{2}\\
&  Q_{RI}Q_{IR}^{t}-Q_{RR}Q_{II}^{t}=-\mathbb{I}_{2}%
\end{align}
If $Q\in Sp\left(  4,\mathbb{R}\right)  \cap O\left(  4,\mathbb{R}\right)  $ then

\begin{itemize}
\item from $Q\in O\left(  4,\mathbb{R}\right)  $%
\begin{align}
&  Q_{RR}Q_{RR}^{t}+Q_{RI}Q_{RI}^{t}=\mathbb{I}_{2}\\
&  Q_{IR}Q_{IR}^{t}+Q_{II}Q_{II}^{t}=\mathbb{I}_{2}\\
&  Q_{IR}Q_{RR}^{t}+Q_{II}Q_{RI}^{t}=\mathbb{O}\Leftrightarrow Q_{IR}%
Q_{RR}^{t}=-Q_{II}Q_{RI}^{t}\\
&  Q_{RR}Q_{IR}^{t}+Q_{RI}Q_{II}^{t}=\mathbb{O}\Leftrightarrow Q_{RR}%
Q_{IR}^{t}=-Q_{RI}Q_{II}^{t}%
\end{align}


\item from $Q\in Sp\left(  4,\mathbb{R}\right)  $%
\begin{align}
&  Q_{RI}Q_{RR}^{t}-Q_{RR}Q_{RI}^{t}=\mathbb{O}\Leftrightarrow Q_{RI}%
Q_{RR}^{t}=Q_{RR}Q_{RI}^{t}\\
&  Q_{II}Q_{IR}^{t}-Q_{IR}Q_{II}^{t}=\mathbb{O}\Leftrightarrow Q_{II}%
Q_{IR}^{t}=Q_{IR}Q_{II}^{t}\\
&  Q_{II}Q_{RR}^{t}-Q_{IR}Q_{RI}^{t}=\mathbb{I}_{2}\\
&  Q_{RI}Q_{IR}^{t}-Q_{RR}Q_{II}^{t}=-\mathbb{I}_{2}%
\end{align}


\item from $Q\in Sp\left(  4,\mathbb{R}\right)  \cap O\left(  4,\mathbb{R}%
\right)  $
\begin{align}
&  Q_{RR}Q_{RR}^{t}+Q_{RI}Q_{RI}^{t}=\mathbb{I}_{2}\,\&\,Q_{II}Q_{RR}%
^{t}-Q_{IR}Q_{RI}^{t}=\mathbb{I}_{2}\\
&  \Rightarrow\mathbb{O}=Q_{RR}Q_{RR}^{t}-Q_{II}Q_{RR}^{t}+Q_{RI}Q_{RI}%
^{t}+Q_{IR}Q_{RI}^{t}\\
&  \Rightarrow Q_{II}Q_{RR}^{t}=\left(  Q_{II}Q_{RR}^{t}\right)  ^{t}\\
&  \Rightarrow Q_{IR}Q_{RI}^{t}=\left(  Q_{IR}Q_{RI}^{t}\right)  ^{t}%
\end{align}%
\begin{align}
&  Q_{RR}Q_{RR}^{t}-Q_{RI}Q_{RI}^{t}=\mathbb{I}_{2}\\
&  Q_{RR}Q_{IR}^{t}-Q_{RI}Q_{II}^{t}=\mathbb{O}\\
&  Q_{IR}Q_{IR}^{t}-Q_{II}Q_{II}^{t}=\mathbb{I}_{2}%
\end{align}

\end{itemize}

If
\begin{equation}
Q=\left(
\begin{array}
[c]{cc}%
Q_{2} & \mathbb{O}\\
\mathbb{O} & Q_{2}%
\end{array}
\right)
\end{equation}
and $Q_{2}\in O\left(  2,\mathbb{R}\right)  $ then
\begin{equation}
Q\in O\left(  2,2,\mathbb{R}\right)  .
\end{equation}
giving
\begin{align}
&  \left(
\begin{array}
[c]{cc}%
Q_{RR} & Q_{RI}\\
Q_{IR} & Q_{II}%
\end{array}
\right)  \left(
\begin{array}
[c]{cc}%
\mathbb{I}_{2} & \mathbb{O}\\
\mathbb{O} & -\mathbb{I}_{2}%
\end{array}
\right)  \left(
\begin{array}
[c]{cc}%
Q_{RR}^{t} & Q_{IR}^{t}\\
Q_{RI}^{t} & Q_{II}^{t}%
\end{array}
\right) \\
&  =\left(
\begin{array}
[c]{cc}%
Q_{RR} & \mathbb{-}Q_{RI}\\
Q_{IR} & -Q_{II}%
\end{array}
\right)  \left(
\begin{array}
[c]{cc}%
Q_{RR}^{t} & Q_{IR}^{t}\\
Q_{RI}^{t} & Q_{II}^{t}%
\end{array}
\right)  =\\
&  \left(
\begin{array}
[c]{cc}%
Q_{RR}Q_{RR}^{t}-Q_{RI}Q_{RI}^{t} & Q_{RR}Q_{IR}^{t}-Q_{RI}Q_{II}^{t}\\
Q_{IR}Q_{RR}^{t}-Q_{II}Q_{RI}^{t} & Q_{IR}Q_{IR}^{t}-Q_{II}Q_{II}^{t}%
\end{array}
\right)  =\left(
\begin{array}
[c]{cc}%
\mathbb{I}_{2} & \mathbb{O}\\
\mathbb{O} & -\mathbb{I}_{2}%
\end{array}
\right) \\
&  Q_{RR}Q_{RR}^{t}-Q_{RI}Q_{RI}^{t}=\mathbb{I}_{2}\\
&  Q_{RR}Q_{IR}^{t}-Q_{RI}Q_{II}^{t}=\mathbb{O}\\
&  Q_{IR}Q_{IR}^{t}-Q_{II}Q_{II}^{t}=\mathbb{I}_{2}%
\end{align}
We have from the precedings that ?
\begin{equation}
Q\in O\left(  4,\mathbb{R}\right)  \cap Sp\left(  4,\mathbb{R}\right)
\Rightarrow Q\in O\left(  2,2,\mathbb{R}\right)
\end{equation}
but
\begin{equation}
Q\in O\left(  2,2,\mathbb{R}\right)  \nRightarrow Q\in O\left(  4,\mathbb{R}%
\right)  \cap Sp\left(  4,\mathbb{R}\right)  .
\end{equation}


\bigskip

The various complex units are given by
\begin{equation}
OJ_{e}O^{t}=O\left\vert e_{3}\right)  \left(  e_{1}\right\vert O^{t}%
+O\left\vert e_{4}\right)  \left(  e_{2}\right\vert O^{t}-O\left\vert
e_{1}\right)  \left(  e_{3}\right\vert O^{t}-O\left\vert e_{2}\right)  \left(
e_{4}\right\vert O^{t}%
\end{equation}


$\lceil$%
\begin{align}
\left(  iJ_{e}\right)  ^{\dagger}  &  =-iJ_{e}^{t}=iJ_{e}\\
\left(  iJ_{e}\right)  ^{\dagger}\left(  iJ_{e}\right)   &  =iJ_{e}iJ_{e}=I
\end{align}%
\begin{equation}
\left(
\begin{array}
[c]{cc}%
0 & -i\\
i & 0
\end{array}
\right)  \left(
\begin{array}
[c]{c}%
a\\
b
\end{array}
\right)  =\lambda\left(
\begin{array}
[c]{c}%
a\\
b
\end{array}
\right)  =\left(
\begin{array}
[c]{c}%
-ib\\
ia
\end{array}
\right)  =\left(
\begin{array}
[c]{c}%
-iia\\
ia
\end{array}
\right)  =\left(
\begin{array}
[c]{c}%
a\\
ia
\end{array}
\right)  =\left(
\begin{array}
[c]{c}%
a\\
ia
\end{array}
\right)
\end{equation}
The solutions for $\lambda$ come from the secular equation
\begin{equation}
\left\vert
\begin{array}
[c]{cc}%
\lambda & -i\\
i & \lambda
\end{array}
\right\vert =\lambda^{2}-1=0\Rightarrow\lambda=\pm1
\end{equation}
which for $\lambda=1$
\begin{equation}
\left(
\begin{array}
[c]{cc}%
0 & -i\\
i & 0
\end{array}
\right)  \left(
\begin{array}
[c]{c}%
a\\
b
\end{array}
\right)  =\left(
\begin{array}
[c]{c}%
-ib\\
ia
\end{array}
\right)  =\left(
\begin{array}
[c]{c}%
-iia\\
ia
\end{array}
\right)  =\left(
\begin{array}
[c]{c}%
a\\
ia
\end{array}
\right)  =\left(
\begin{array}
[c]{c}%
a\\
ia
\end{array}
\right)
\end{equation}
gives the eigen vector
\begin{equation}
\left(
\begin{array}
[c]{c}%
a\\
b
\end{array}
\right)  =\frac{1}{\sqrt{2}}\left(
\begin{array}
[c]{c}%
1\\
i
\end{array}
\right)
\end{equation}
and for $\lambda=-1$
\begin{equation}
\left(
\begin{array}
[c]{cc}%
0 & -i\\
i & 0
\end{array}
\right)  \left(
\begin{array}
[c]{c}%
a\\
b
\end{array}
\right)  =\left(
\begin{array}
[c]{c}%
ib\\
-ia
\end{array}
\right)  =\left(
\begin{array}
[c]{c}%
-iia\\
ia
\end{array}
\right)  =\left(
\begin{array}
[c]{c}%
a\\
ia
\end{array}
\right)  =-\left(
\begin{array}
[c]{c}%
a\\
ia
\end{array}
\right)
\end{equation}
gives the eigen vector
\begin{equation}
\left(
\begin{array}
[c]{c}%
a\\
b
\end{array}
\right)  =\frac{1}{\sqrt{2}}\left(
\begin{array}
[c]{c}%
1\\
-i
\end{array}
\right)  ,
\end{equation}
which is displayed by
\begin{align*}
\left(
\begin{array}
[c]{cc}%
0 & -i\\
i & 0
\end{array}
\right)  \frac{1}{\sqrt{2}}\left(
\begin{array}
[c]{c}%
1\\
i
\end{array}
\right)   &  =\frac{1}{\sqrt{2}}\left(
\begin{array}
[c]{c}%
1\\
i
\end{array}
\right) \\
\left(
\begin{array}
[c]{cc}%
0 & -i\\
i & 0
\end{array}
\right)  \frac{1}{\sqrt{2}}\left(
\begin{array}
[c]{c}%
1\\
-i
\end{array}
\right)   &  =\frac{1}{\sqrt{2}}\left(
\begin{array}
[c]{c}%
-1\\
i
\end{array}
\right)  =-\frac{1}{\sqrt{2}}\left(
\begin{array}
[c]{c}%
1\\
-i
\end{array}
\right)  .
\end{align*}


This is confirmed by expanding
\begin{equation}
iJ_{e}=\left\vert v_{1}\right\rangle \left\langle v_{1}\right\vert -\left\vert
v_{2}\right\rangle \left\langle v_{2}\right\vert
\end{equation}
to obtain%
\begin{equation}
\left(
\begin{array}
[c]{cc}%
0 & -i\\
i & 0
\end{array}
\right)  =\frac{1}{2}\left(
\begin{array}
[c]{c}%
1\\
i
\end{array}
\right)  \left(
\begin{array}
[c]{cc}%
1 & -i
\end{array}
\right)  -\frac{1}{2}\left(
\begin{array}
[c]{c}%
1\\
-i
\end{array}
\right)  \left(
\begin{array}
[c]{cc}%
1 & i
\end{array}
\right)  =\frac{1}{2}\left(
\begin{array}
[c]{cc}%
1 & -i\\
i & 1
\end{array}
\right)  -\frac{1}{2}\left(
\begin{array}
[c]{cc}%
1 & i\\
-i & 1
\end{array}
\right)  =\left(
\begin{array}
[c]{cc}%
0 & i\\
-i & 0
\end{array}
\right)
\end{equation}
corresponding to
\begin{equation}
J_{e}=-i\left\{  \left\vert v_{1}\right\rangle \left\langle v_{1}\right\vert
-\left\vert v_{2}\right\rangle \left\langle v_{2}\right\vert \right\}
\end{equation}
hence
\begin{equation}
-i\left(
\begin{array}
[c]{cc}%
0 & -i\\
i & 0
\end{array}
\right)  =\frac{1}{2}\left(
\begin{array}
[c]{c}%
-i\\
1
\end{array}
\right)  \left(
\begin{array}
[c]{cc}%
1 & -i
\end{array}
\right)  +\frac{1}{2}\left(
\begin{array}
[c]{c}%
i\\
1
\end{array}
\right)  \left(
\begin{array}
[c]{cc}%
1 & i
\end{array}
\right)
\end{equation}
corresponding to%
\begin{equation}
iJ_{4e}=\left\vert v_{1}\right\rangle \left\langle v_{1}\right\vert
+\left\vert v_{2}\right\rangle \left\langle v_{2}\right\vert -\left(
\left\vert v_{3}\right\rangle \left\langle v_{3}\right\vert +\left\vert
v_{4}\right\rangle \left\langle v_{4}\right\vert \right)  =P_{1}-P_{2}%
\end{equation}


$\rfloor$

Define
\begin{align}
P_{R}  &  =\left\vert e_{1}\right)  \left(  e_{1}\right\vert +\left\vert
e_{2}\right)  \left(  e_{2}\right\vert \\
P_{I}  &  =\left\vert e_{3}\right)  \left(  e_{3}\right\vert +\left\vert
e_{4}\right)  \left(  e_{4}\right\vert
\end{align}
then
\begin{equation}
J_{4e}=\left\vert e_{3}\right)  \left(  e_{1}\right\vert +\left\vert
e_{4}\right)  \left(  e_{2}\right\vert -\left\vert e_{1}\right)  \left(
e_{3}\right\vert -\left\vert e_{2}\right)  \left(  e_{4}\right\vert
\end{equation}
and%
\begin{align}
P_{R}J_{4e}P_{I}  &  =-\left\vert e_{1}\right)  \left(  e_{3}\right\vert
-\left\vert e_{2}\right)  \left(  e_{4}\right\vert \\
P_{I}J_{4e}P_{R}  &  =\left\vert e_{3}\right)  \left(  e_{1}\right\vert
+\left\vert e_{4}\right)  \left(  e_{2}\right\vert .
\end{align}


From the preceding $Q\in Sp\left(  4,\mathbb{R}\right)  \cap O\left(
4,\mathbb{R}\right)  $ if
\begin{equation}
\left[  Q,P_{R}\right]  =\left[  Q,P_{I}\right]  =0
\end{equation}


\section{Complexification as Representations of $\mathbb{C}$}

Consider a real even dimensional Hilbert space $\mathcal{V}$
\begin{equation}
\dim\left\{  \mathcal{V}\right\}  =2s
\end{equation}
and the set of bounded linear transformations $\mathcal{B}\left(
\mathcal{V}\right)  .$

An orthogonal complexification is given by a complex unit $J_{e}\in
\mathcal{B}\left(  \mathcal{V}\right)  $ that has the properties
\begin{gather}
J_{e}^{t}=-J_{e}\,\&J_{e}^{2}=-I_{V}\\
\Downarrow\\
J_{e}\in O\left(  \mathbb{R}^{4}\right)
\end{gather}
as%
\begin{align}
J_{e}^{2}  &  =-I_{\mathcal{V}}\\
J_{e}^{t}  &  =-J_{e}%
\end{align}
Where
\begin{equation}
J_{e}=\sum\limits_{1\leq j\leq s}\left\{  \left\vert e_{j}\right)  \left(
e_{j+s}\right\vert -\left\vert e_{j+s}\right)  \left(  e_{j}\right\vert
\right\}  .
\end{equation}
The identity belonging to $\mathcal{B}\left(  \mathcal{V}\right)  $ can be
expressed as
\begin{equation}
I_{\mathcal{V}}=\sum\limits_{1\leq j\leq2s}\left\vert e_{j}\right)  \left(
e_{j}\right\vert =P_{e}+P_{\bar{e}},
\end{equation}
where
\begin{align}
P_{e}  &  =\sum\limits_{1\leq j\leq s}\left\vert e_{j}\right)  \left(
e_{j}\right\vert \\
P_{\bar{e}}  &  =\sum\limits_{s+1\leq j\leq2s}\left\vert e_{j}\right)  \left(
e_{j}\right\vert .
\end{align}


The complex numbers $\mathbb{C}$ can be represented in $\mathcal{B}\left(
\mathcal{V}\right)  $ by
\begin{equation}
K_{e}\left(  z\right)  =\operatorname{Re}zI_{e}+\operatorname{Im}zJ_{e}%
=I_{e}\operatorname{Re}z+J_{e}\operatorname{Im}z;z\in\mathbb{C}%
\end{equation}
$\lceil$%
\begin{equation}
z_{1}z_{2}=\left(  \operatorname{Re}z_{1}+i\operatorname{Im}z_{1}\right)
\left(  \operatorname{Re}z_{2}+i\operatorname{Im}z_{2}\right)  =\left(
\operatorname{Re}z_{1}\operatorname{Re}z_{2}-\operatorname{Im}z_{1}%
\operatorname{Im}z_{2}\right)  +i\left(  \operatorname{Im}z_{1}%
\operatorname{Re}z_{2}+\operatorname{Re}z_{1}\operatorname{Im}z_{2}\right)
\end{equation}%
\begin{align}
K_{e}\left(  z_{1}z_{2}\right)   &  =\left(  \operatorname{Re}z_{1}%
\operatorname{Re}z_{2}-\operatorname{Im}z_{1}\operatorname{Im}z_{2}\right)
I_{e}+\left(  \operatorname{Im}z_{1}\operatorname{Re}z_{2}+\operatorname{Re}%
z_{1}\operatorname{Im}z_{2}\right)  J_{e}\\
K_{e}\left(  z_{1}\right)  K_{e}\left(  z_{2}\right)   &  =\left(
\operatorname{Re}z_{1}I_{e}+\operatorname{Im}z_{1}J_{e}\right)  \left(
\operatorname{Re}z_{2}I_{e}+\operatorname{Im}z_{2}J_{e}\right)  =\left(
\operatorname{Re}z_{1}\operatorname{Re}z_{2}-\operatorname{Im}z_{1}%
\operatorname{Im}z_{2}\right)  I_{e}+\left(  \operatorname{Im}z_{1}%
\operatorname{Re}z_{2}+\operatorname{Re}z_{1}\operatorname{Im}z_{2}\right)
J_{e}\\
K_{e}\left(  z_{1}z_{2}\right)   &  =K_{e}\left(  z_{1}\right)  K_{e}\left(
z_{2}\right)
\end{align}


$\rfloor$

inducing the complexification
\begin{align}
\left\vert a\right)   &  =\left\vert e_{1},e_{2},e_{3},e_{4}\right)  \left(
\begin{array}
[c]{c}%
a_{11}\\
a_{21}\\
a_{12}\\
a_{22}%
\end{array}
\right)  =\left\vert e_{1},e_{2},J_{e}e_{1},J_{e}e_{2}\right)  \left(
\begin{array}
[c]{c}%
a_{11}\\
a_{21}\\
a_{12}\\
a_{22}%
\end{array}
\right) \\
&  =\left\vert e_{1}\right)  a_{11}+\left\vert e_{2}\right)  a_{21}%
+J_{e}\left\vert e_{1}\right)  a_{12}+J_{e}\left\vert e_{2}\right)  a_{22}\\
&  =a_{11}I_{2s}\left\vert e_{1}\right)  +a_{21}I_{2s}\left\vert e_{2}\right)
+a_{12}J_{e}\left\vert e_{1}\right)  +a_{22}J_{e}\left\vert e_{2}\right) \\
&  =\left(  a_{11}I_{2s}+a_{12}J_{e}\right)  \left\vert e_{1}\right)  +\left(
a_{21}I_{2s}+a_{22}J_{e}\right)  \left\vert e_{2}\right) \\
&  =K_{e}\left(  z_{1}\right)  \left\vert e_{1}\right)  +K_{e}\left(
z_{2}\right)  \left\vert e_{2}\right) \\
&  =\left\vert K_{e}\left(  z_{1}\right)  ,K_{e}\left(  z_{2}\right)  \right)
\left(
\begin{array}
[c]{c}%
\left\vert e_{1}\right) \\
\left\vert e_{2}\right)
\end{array}
\right)  ,
\end{align}
where%
\begin{align}
K_{e}\left(  z_{1}\right)   &  =a_{11}I_{2s}+a_{12}J_{e}=\operatorname{Re}%
z_{1}I_{e}+\operatorname{Im}z_{1}J_{e}=\left\vert I_{2s},J_{e}\right)  \left(
%
\begin{array}
[c]{c}%
\operatorname{Re}z_{1}\\
\operatorname{Im}z_{1}%
\end{array}
\right) \\
K_{e}\left(  z_{2}\right)   &  =a_{21}I_{2s}+a_{22}J_{e}=\operatorname{Re}%
z_{2}I_{e}+\operatorname{Im}z_{2}J_{e}=\left\vert I_{2s},J_{e}\right)  \left(
%
\begin{array}
[c]{c}%
\operatorname{Re}z_{2}\\
\operatorname{Im}z_{2}%
\end{array}
\right)  .
\end{align}
In order to see that this a complexification one expands
\begin{align}
K_{e}\left(  z_{1}\right)  \left\vert e_{1}\right)  +K_{e}\left(
z_{2}\right)  \left\vert e_{2}\right)   &  =\left\vert I_{2s},J_{e}\right)
\left(
\begin{array}
[c]{c}%
\operatorname{Re}z_{1}\\
\operatorname{Im}z_{1}%
\end{array}
\right)  \left\vert e_{1}\right)  +\left\vert I_{2s},J_{e}\right)  \left(
\begin{array}
[c]{c}%
\operatorname{Re}z_{2}\\
\operatorname{Im}z_{2}%
\end{array}
\right)  \left\vert e_{2}\right) \\
&  =\left\vert I_{2s},J_{e}\right)  \left\{  \left(
\begin{array}
[c]{c}%
\operatorname{Re}z_{1}\\
\operatorname{Im}z_{1}%
\end{array}
\right)  \left\vert e_{1}\right)  +\left(
\begin{array}
[c]{c}%
\operatorname{Re}z_{2}\\
\operatorname{Im}z_{2}%
\end{array}
\right)  \left\vert e_{2}\right)  \right\} \\
&  =\left\vert I_{2s},J_{e}\right)  \left\{  \left\vert \left(
\begin{array}
[c]{c}%
\operatorname{Re}z_{1}\\
\operatorname{Im}z_{1}%
\end{array}
\right)  ,\left(
\begin{array}
[c]{c}%
\operatorname{Re}z_{2}\\
\operatorname{Im}z_{2}%
\end{array}
\right)  \right)  \left(
\begin{array}
[c]{c}%
\left\vert e_{1}\right) \\
\left\vert e_{2}\right)
\end{array}
\right)  \right\} \\
&  =\left\vert I_{2s},J_{e}\right)  \left\{  \left\vert z_{1},z_{2}\right)
\left(
\begin{array}
[c]{c}%
\left\vert e_{1}\right) \\
\left\vert e_{2}\right)
\end{array}
\right)  \right\}  .
\end{align}
Also
\begin{align}
K_{e}\left(  z_{1}\right)  \left\vert e_{1}\right)  +K_{e}\left(
z_{2}\right)  \left\vert e_{2}\right)   &  =\left\vert I_{2s},J_{e}\right)
\left\{  \left\vert \left\vert e_{1}\right)  ,\left\vert e_{2}\right)
\right)  \left(
\begin{array}
[c]{c}%
\left(
\begin{array}
[c]{c}%
\operatorname{Re}z_{1}\\
\operatorname{Im}z_{1}%
\end{array}
\right) \\
\left(
\begin{array}
[c]{c}%
\operatorname{Re}z_{2}\\
\operatorname{Im}z_{2}%
\end{array}
\right)
\end{array}
\right)  \right\} \\
&  =\left\vert I_{2s},J_{e}\right)  \left\{  \left\vert \left\vert
e_{1}\right)  ,\left\vert e_{2}\right)  \right)  \left(
\begin{array}
[c]{c}%
z_{1}\\
z_{2}%
\end{array}
\right)  \right\}  ,
\end{align}
where
\begin{align}
z  &  \leftrightarrow\left(
\begin{array}
[c]{c}%
\operatorname{Re}z\\
\operatorname{Im}z
\end{array}
\right)  \leftrightarrow\left(  1,i\right)  \left(
\begin{array}
[c]{c}%
\operatorname{Re}z\\
\operatorname{Im}z
\end{array}
\right)  =1\left(  \operatorname{Re}z\right)  +i\left(  \operatorname{Im}%
z\right)  \in\mathcal{B}\left(  \mathbb{R}^{2}\right) \\
1,i  &  \in\mathcal{B}\left(  \mathbb{R}^{2}\right)  .
\end{align}


\section{Representation of $\mathcal{V}$ in $\mathcal{B}\left(  \mathcal{V}%
\right)  $}

The space $\mathcal{V}$ can be represented in $\mathcal{B}\left(
\mathcal{V}\right)  $ by the linear injection
\begin{align}
K_{e}  &  :\mathcal{V\hookrightarrow B}\left(  \mathcal{V}\right) \\
\left\vert e_{j}\right)   &  \rightarrow\left\vert e_{j}\right)  \left(
e_{j}\right\vert \\
K_{e}\left(  a\right)   &  =\sum\limits_{1\leq j\leq2s}\left\vert
e_{j}\right)  \left(  e_{j}\right\vert R_{e}\left(  a\right)  _{j}\\
\left\vert a\right)   &  =\sum\limits_{1\leq j\leq2s}\left\vert e_{j}\right)
R_{e}\left(  a\right)  _{j}.
\end{align}
$K_{e}$ is a Hilbert space isomorphism as
\begin{align}
K_{e}  &  :\mathcal{V\hookrightarrow}K_{e}\left(  \mathcal{V}\right)
\subset\mathcal{B}\left(  \mathcal{V}\right) \\
\left(  \left.  a\right\vert b\right)   &  =\operatorname{Tr}\left\{
K_{e}\left(  a\right)  ^{t}K_{e}\left(  b\right)  \right\}  .
\end{align}%
\begin{equation}
K_{e}\left(  \mathcal{V}_{\mathbb{C}}\right)  =\left\{  \left(  \left\vert
e_{1}\right)  \left(  e_{1}\right\vert ,\cdots,\left\vert e_{2s}\right)
\left(  e_{2s}\right\vert \right)  \left(
\begin{array}
[c]{c}%
a_{1}\\
\vdots\\
a_{2s}%
\end{array}
\right)  \right\}
\end{equation}%
\begin{equation}
=\left\{  \left(  \left\vert e_{1}\right)  \left(  e_{1}\right\vert
,\cdots,\left\vert e_{s}\right)  \left(  e_{s}\right\vert ,\left\vert
e_{s+1}\right)  \left(  e_{s+1}\right\vert ,\cdots,\left\vert e_{2s}\right)
\left(  e_{2s}\right\vert \right)  \left(
\begin{array}
[c]{c}%
a_{1}\\
\vdots\\
a_{s}\\
a_{s+1}\\
\vdots\\
a_{2s}%
\end{array}
\right)  \right\}
\end{equation}
%

\begin{align}
=  &  \left(  \left(  \left\vert e_{1}\right)  \left(  e_{1}\right\vert
,\cdots,\left\vert e_{s}\right)  \left(  e_{s}\right\vert \right)  \right)
\left(  \left(
\begin{array}
[c]{c}%
a_{1}\\
\vdots\\
a_{s}%
\end{array}
\right)  \right) \\
&  +\left(  \left(  J_{e}\left\vert e_{1}\right)  \left(  e_{1}\right\vert
J_{e}^{t},\cdots,J_{e}\left\vert e_{s}\right)  \left(  e_{s}\right\vert
J_{e}^{t}\right)  \right)  \left(  \left(
\begin{array}
[c]{c}%
a_{s+1}\\
\vdots\\
a_{2s}%
\end{array}
\right)  \right)
\end{align}%
\begin{equation}
=\left\{  \left(  \left\vert e_{1}\right)  \left(  e_{1}\right\vert
,\cdots,\left\vert e_{s}\right)  \left(  e_{s}\right\vert ,J_{e}\left\vert
e_{1}\right)  \left(  e_{1}\right\vert J_{e}^{t},\cdots,J_{e}\left\vert
e_{s}\right)  \left(  e_{s}\right\vert J_{e}^{t}\right)  \left(
\begin{array}
[c]{c}%
a_{1}\\
\vdots\\
a_{s}\\
a_{s+1}\\
\vdots\\
a_{2s}%
\end{array}
\right)  \right\}
\end{equation}%
\begin{equation}
=\left\{  \left(  \left(  \left\vert e_{1}\right)  \left(  e_{1}\right\vert
,\cdots,\left\vert e_{s}\right)  \left(  e_{s}\right\vert \right)  ,\left(
J_{e}\left\vert e_{1}\right)  \left(  e_{1}\right\vert J_{e}^{t},\cdots
,J_{e}\left\vert e_{s}\right)  \left(  e_{s}\right\vert J_{e}^{t}\right)
\right)  \left(
\begin{array}
[c]{c}%
\left(
\begin{array}
[c]{c}%
a_{1}\\
\vdots\\
a_{s}%
\end{array}
\right) \\
\left(
\begin{array}
[c]{c}%
a_{s+1}\\
\vdots\\
a_{2s}%
\end{array}
\right)
\end{array}
\right)  \right\}
\end{equation}%
\begin{equation}
=\left\{  \left(  \left(  \left\vert e_{1}\right)  \left(  e_{1}\right\vert
,\cdots,\left\vert e_{s}\right)  \left(  e_{s}\right\vert \right)  ,\left(
\widehat{J_{e}}\left(  \left\vert e_{1}\right)  \left(  e_{1}\right\vert
\right)  ,\cdots,\widehat{J_{e}}\left(  \left\vert e_{s}\right)  \left(
e_{s}\right\vert \right)  \right)  \right)  \left(
\begin{array}
[c]{c}%
\left(
\begin{array}
[c]{c}%
a_{1}\\
\vdots\\
a_{s}%
\end{array}
\right) \\
\left(
\begin{array}
[c]{c}%
a_{s+1}\\
\vdots\\
a_{2s}%
\end{array}
\right)
\end{array}
\right)  \right\}  .
\end{equation}%
\begin{align}
J_{e}\left\vert e_{k}\right)  \left(  e_{k}\right\vert J_{e}^{t}  &
=\left\vert e_{k+s}\right)  \left(  e_{k+s}\right\vert \\
J_{e}\left\vert e_{k+s}\right)  \left(  e_{k+s}\right\vert J_{e}^{t}  &
=\left\vert e_{k}\right)  \left(  e_{k}\right\vert
\end{align}


\bigskip

\bigskip

\bigskip%

\begin{align}
\left(
\begin{array}
[c]{cc}%
Q_{RR} & Q_{RI}\\
Q_{IR} & Q_{II}%
\end{array}
\right)  \left(
\begin{array}
[c]{cc}%
\mathbb{I}_{2} & \mathbb{O}\\
\mathbb{O} & -\mathbb{I}_{2}%
\end{array}
\right)   &  =\left(
\begin{array}
[c]{cc}%
\mathbb{I}_{2} & \mathbb{O}\\
\mathbb{O} & -\mathbb{I}_{2}%
\end{array}
\right)  \left(
\begin{array}
[c]{cc}%
Q_{RR} & Q_{RI}\\
Q_{IR} & Q_{II}%
\end{array}
\right)  ^{t}\\
&  \Updownarrow\\
\left(
\begin{array}
[c]{cc}%
Q_{RR} & -Q_{RI}\\
Q_{IR} & -Q_{II}%
\end{array}
\right)   &  =\left(
\begin{array}
[c]{cc}%
Q_{RR}^{t} & -Q_{IR}^{t}\\
Q_{RI}^{t} & -Q_{II}^{t}%
\end{array}
\right) \\
&  \Updownarrow\\
\left(
\begin{array}
[c]{cc}%
Q_{RR} & Q_{RI}\\
Q_{IR} & Q_{II}%
\end{array}
\right)   &  =\left(
\begin{array}
[c]{cc}%
Q_{RR} & \mathbb{O}\\
\mathbb{O} & Q_{II}%
\end{array}
\right)
\end{align}%
\begin{align}
\left(
\begin{array}
[c]{cc}%
Q_{RR} & \mathbb{O}\\
\mathbb{O} & Q_{II}%
\end{array}
\right)  \left(
\begin{array}
[c]{cc}%
\mathbb{O} & \mathbb{-I}_{2}\\
\mathbb{I}_{2} & \mathbb{O}%
\end{array}
\right)  \left(
\begin{array}
[c]{cc}%
Q_{RR} & \mathbb{O}\\
\mathbb{O} & Q_{II}%
\end{array}
\right)  ^{t}  &  =\left(
\begin{array}
[c]{cc}%
\mathbb{O} & \mathbb{-I}_{2}\\
\mathbb{I}_{2} & \mathbb{O}%
\end{array}
\right) \\
\left(
\begin{array}
[c]{cc}%
\mathbb{O} & \mathbb{-}Q_{RR}\\
Q_{II} & \mathbb{O}%
\end{array}
\right)  \left(
\begin{array}
[c]{cc}%
Q_{RR}^{t} & \mathbb{O}\\
\mathbb{O} & Q_{II}^{t}%
\end{array}
\right)   &  =\left(
\begin{array}
[c]{cc}%
\mathbb{O} & \mathbb{-}Q_{RR}Q_{II}^{t}\\
Q_{II}Q_{RR}^{t} & \mathbb{O}%
\end{array}
\right)  =\left(
\begin{array}
[c]{cc}%
\mathbb{O} & \mathbb{-I}_{2}\\
\mathbb{I}_{2} & \mathbb{O}%
\end{array}
\right)
\end{align}



\end{document}